\documentclass[12pt,a4paper]{article}
\usepackage[T1]{fontenc}
\usepackage[utf8]{inputenc}
\usepackage{amsmath,amssymb,amsfonts}
\usepackage{geometry}
\usepackage{enumitem}
\usepackage{booktabs}
\usepackage{longtable}
\usepackage{array}
\usepackage[hidelinks]{hyperref}

\geometry{margin=2.5cm}
\setlength{\parskip}{0.6ex}
\setlength{\parindent}{0pt}
\emergencystretch=2em
\widowpenalty=1000
\clubpenalty=1000

\newcommand{\id}[1]{\texttt{#1}}
\newcommand{\mech}[1]{\paragraph{Mechanism.} #1}
\newcommand{\why}[1]{\paragraph{Why this design.} #1}
\newcommand{\alt}[1]{\paragraph{Alternatives.} #1}

\title{Model Logic: Design Choices, Justifications and Alternatives\\[0.4em]
\large Cancer Screening Microsimulation Model (MedicalModel2024)}
\author{Model documentation}
\date{\today}

\begin{document}
\maketitle
\tableofcontents
\newpage

\section{Purpose, Scope and How to Read This Document}

This document describes \emph{what the model does and why it is built that way}. It is deliberately
neither a code manual nor a port review. For almost every element of the model there is more than one
defensible way to do it, so the useful information is not only the choice that was made but the
alternatives that were rejected and the reason for rejecting them --- because those reasons determine
where the model is strong, and where a different dataset or a different question would force a
different design.

Every modelling element is presented with the same three-part template:

\begin{enumerate}[nosep]
  \item \textbf{Mechanism} --- the rule the model actually applies, stated so that it could be
        re-implemented from this description alone.
  \item \textbf{Why this design} --- what the choice buys: identifiability from the available data,
        epidemiological plausibility, interpretability, or computational cost.
  \item \textbf{Alternatives} --- other defensible designs, why each was not used here, and what would
        have to change to use it.
\end{enumerate}

Two sibling documents cover the other angles and are not repeated here:

\begin{itemize}[nosep]
  \item \id{model\_for\_dummies.pdf} --- a plain-language walkthrough of the model and the web interface.
  \item \id{python\_port\_analysis.pdf} --- the technical and statistical review: state variables,
        statistical corrections, performance, and methodological improvements.
\end{itemize}

\textbf{Reference configuration.} Every number quoted below as an illustration comes from one
configuration: the GLOBOCAN colorectal datasets (USA aggregate and staging files), $N = 20\,000$ agents,
seed 2022, 30 simulated years, screening at ages 50--60 every second year with 50\% participation. That
run yields 1\,142 cancer cases, 200 patients cured clinically, 5 cured after screen detection, 12\,334
survivors at the end, and a total cost of 13\,509 units. Numbers from other runs are not comparable
one-to-one, because the model is stochastic.

\textbf{A note on honesty.} Where an element exists in the code but has no effect on the results, it is
listed explicitly in Section~\ref{sec:inactive} rather than silently omitted. Where the model
\emph{cannot} determine a quantity from the available data, that is stated at the point where it
matters.

\section{Model Class, Time and Population Frame}

\subsection{Why an agent-based microsimulation}

\mech{The model creates $N$ independent agents. Each agent is given one complete cancer history ---
whether and when a first malignant cell appears, how fast the tumour grows, what stage it has reached by
the time of clinical diagnosis, whether treatment cures it, and at what age it dies --- and is then
followed year by year. Nothing is aggregated into fractions; the population statistics are obtained by
counting agents after the fact.}

\why{Four properties of the question force this design.

\emph{(i) The answer is a difference of tails, not a mean.} Screening is justified by a small number of
lives shifted out of the cancer-death column, and by the years of life those people gain. Both are
properties of individual counterfactuals, and they are exactly the quantities a model of averages cannot
represent: the same individual must be observable with and without the intervention. The model therefore
stores \emph{both} death ages for every patient (untreated and screening-affected), so the survivors-per-year
curve is computed from within-individual differences rather than from two independent populations.
Heterogeneity and the tail go together: the people who benefit are not the average person.

\emph{(ii) Screening is history-dependent.} Who is screened at age 58 depends on participation at 56, on
whether a cancer was already found, and on the age window. In a state-based model this forces the state
space to carry the whole screening history. In a microsimulation it is a few boolean flags per agent.

\emph{(iii) The heterogeneity is structural, not decorative.} Individual growth rate
($B_{\text{ind}}$), per-stage aggressiveness, age-dependent cure efficiency and per-stage treatment
efficiency all act multiplicatively on the outcome. Representing any of them faithfully in a cohort model
means multiplying the number of states by the number of levels; combining them is combinatorially
explosive.

\emph{(iv) The model must reproduce a stage mix, not a mean stage.} The stage at diagnosis is a
consequence of how long the tumour has been growing and how fast, so it is inherently an individual-level
crossing event (Section~\ref{sec:stage}).}

\alt{\begin{itemize}[nosep]
  \item \textbf{Cohort (deterministic) Markov / multi-state life table.} The classical alternative, used
        for many cancer screening analyses because it is fast, deterministic and easy to communicate.
        Rejected because it cannot carry the per-individual counterfactual that the ``lives saved''
        output is defined by, and because the state space grows multiplicatively with heterogeneity.
  \item \textbf{Age-structured PDE (McKendrick--von Foerster).} Analytically tractable and fast for
        mean-field questions. Rejected for the same reason, plus a practical one: screening makes the
        transfer rates time- and history-dependent, which turns the PDE into a large system.
  \item \textbf{Renewal equation / integral projection model.} Very efficient for incidence and stage
        shift. Rejected because cause-specific mortality and the within-individual trade of cancer death
        against other-cause death are awkward to express.
  \item \textbf{Semi-Markov with sojourn-time distributions.} A good compromise and a genuine option.
        Rejected here because the stage-crossing mechanism is already a semi-Markov process with an
        explicit growth model; adding a second latent progression structure would double-count.
  \item \textbf{Compartmental ODE / system dynamics.} Rejected: no individual variability, no screening
        history, and ``who benefits'' becomes unanswerable.
  \item \textbf{Decision tree or Bayesian network.} Appropriate for a single one-off decision, not for a
        lifetime intervention with competing risks.
\end{itemize}}

\subsection{Why one-year discrete time steps}

\mech{Time is the integer year counter \id{current\_date}. Ages are integers
($\text{age} = \id{current\_date} - \id{date\_birth}$). At each step the model (a) removes the agents
whose death age equals their current age, (b) runs a screening round if the year is on the screening
grid, and (c) updates the statistics. All events within a year are resolved at year boundaries.}

\why{The time step is chosen to match the granularity of the calibration data and of the intervention
being modelled, not for convenience. The aggregate data are annual age-banded tables; the programme being
evaluated is defined on an annual grid (``invite at ages 50--60 every two years''); and the natural-history
processes involved are slow compared with a year, with one exception noted below. Annual integer ages make
three internal devices exact rather than approximate: the death comparison is an integer equality, the
stage at screen detection is a lookup of the stored crossing age against the current age, and the
exported agent table is directly readable.}

\why{The price is explicit: within a year, events are simultaneous and unordered, no event can occur
twice in the same year, and a process faster than a year --- the terminal phase after the last stage
crossing --- is discretised. That last point is not a footnote here: the fitted cancer-death hazard gives
a mean terminal survival of roughly three years (Section~\ref{sec:calib}), so a one-year grid is adequate
but not generous.}

\alt{\begin{itemize}[nosep]
  \item \textbf{Continuous time with exponential clocks (Gillespie-style next-event loop).} Exact
        within-year dynamics and exact screening instants; requires an exponential assumption for every
        transition and a much larger implementation surface. This is the most defensible upgrade if the
        timing of the terminal phase ever becomes the object of study.
  \item \textbf{Monthly or quarterly steps.} Finer, but multiplies cost by 12 or 4 while the calibration
        data cannot support sub-annual targets, so the extra resolution would not be controlled by data.
  \item \textbf{Event-driven simulation with a priority queue.} Natural for continuous time; more
        machinery than the annual grid needs.
  \item \textbf{Hybrid: annual screening grid with continuous natural history.} A reasonable middle
        ground, but it complicates reproducibility and the exported agent table for no gain at the
        current level of data support.
\end{itemize}}

\subsection{Why a closed cohort and a 30-year horizon}

\mech{There are no births and no migration: the population ages and dies. The initial age structure is
taken from the observed population, so the cohort already contains all ages at $t=0$. The horizon is
\id{years\_to\_simulate} (30 by default).}

\why{A closed cohort isolates the effect being estimated. If births and migration were simulated, the
observed mortality difference between the screening and no-screening scenarios would mix the screening
effect with demographic drift, and the model would need fertility and migration assumptions it cannot
validate against anything. Because the initial age distribution is cross-sectional --- people enter at
their current age --- the cohort is already realistic at $t=0$, so no burn-in is needed and the first
year of output is usable; this matters for a screening window that opens at age 50. Thirty years covers
the whole screening window twice over and allows the mortality difference to accumulate, while keeping
the run interactive: the reference 20\,000-agent, 30-year run completes in a couple of seconds.}

\alt{\begin{itemize}[nosep]
  \item \textbf{Birth cohort followed from age 0.} Cleanest for lifetime outcomes and for interactions
        with childhood exposures; requires a burn-in of a century and many more agents to reach the same
        statistical power in the ages 50--80 where the decisions are.
  \item \textbf{Open population with demographic projection.} More realistic for long-run forecasts;
        introduces fertility and migration assumptions this model has no data to support, and
        contaminates the screening contrast.
  \item \textbf{Steady-state (stationary) population.} Attractive if only long-run rates matter; adds a
        demographic-equilibrium assumption without helping the screening question.
  \item \textbf{Hybrid entry at mid-life.} Entering a cohort at age 40--50 would shorten runs further, but
        discards the interaction between the age at which screening starts and the pre-screening history.
\end{itemize}}

\section{Demography: Initial Age Structure and Competing Mortality}

\subsection{Initial age structure}

\mech{Every agent is assigned a starting age by sampling from the observed population age structure: the
aggregate data file's \id{population} column, one value per year of age, is normalised and used as a
discrete distribution; ages are drawn by the inverse-CDF method (a uniform draw is located in the
cumulative array). The agent's birth year is then set to minus that age, so absolute age is simply
``time since birth''.}

\why{The cross-sectional age structure is the correct initial condition for a closed cohort: it fixes both
the number of people at risk and the age profile of that risk at the moment the simulation starts. Using
the file as a probability mass rather than as absolute counts means the file's units cancel, so the model
does not depend on whether the column is in persons, thousands or percentages. Single years of age are
used because the data are already available at that resolution; the model imposes no binning of its own,
so any loss of detail is the data's choice rather than the model's.}

\alt{\begin{itemize}[nosep]
  \item \textbf{A uniform or single-peaked age distribution.} Rejected: the age distribution is the
        dominant driver of both incidence and mortality, so an unrealistic one would corrupt every
        output.
  \item \textbf{Age bands with a step function.} Cheaper and closer to how source data are often
        published; it loses intra-band resolution, which matters near the screening window.
  \item \textbf{A fitted parametric age profile (Gompertz, Weibull, or a spline).} Smooths sampling noise
        and extrapolates, but imposes a shape the data do not require, and any misfit propagates into
        every rate.
  \item \textbf{Empirical microdata (a census or survey sample of individual ages).} The most faithful
        option; it needs data the packaged model does not ship with, and the aggregate route already
        reproduces the marginal distribution.
  \item \textbf{Starting everyone at the same age (e.g.\ all at 50).} Rejected: it destroys the age
        composition, removes the youngest and oldest, and makes the screening window a degenerate cohort
        boundary.
\end{itemize}}

\subsection{Non-cancer mortality as a competing risk}

\mech{Each agent also receives a \emph{natural death age}, sampled in the same way from a distribution
built from the observed annual all-cause death counts (\id{deaths all} divided by the total population,
normalised). Agents whose sampled death age is not greater than their starting age are resampled once; if
the redraw still fails the death age is set to starting age $+1$, so the agent dies in the first simulated
year. Thereafter the natural death age is fixed, and everything the model does is compared against it.}

\why{Addressing non-cancer mortality as a single fixed individual ``death age'', drawn once, buys three
things.

\emph{(i) A competing risk becomes an arithmetic comparison.} With one natural death age and one cancer
death age per agent, everything that depends on their interaction --- did the cancer or something else kill
this person first, was this person cured, did screening save them, how many life-years were gained --- is a
comparison of two numbers. This is what makes the paired screening/no-screening accounting exact rather
than approximate: there is no control population to re-simulate and no extra sampling noise in the
contrast.

\emph{(ii) It is distributionally equivalent to a survival model without the bookkeeping.} Drawing the
death age from the life table is marginally the same as removing an annual fraction of survivors, but it
replaces a running state variable and an iterative update with one stored integer.

\emph{(iii) The conditional redraw is the right correction.} Because agents enter at their current age,
their death age must be conditional on already being alive; resampling the minority whose first draw falls
below the entry age implements exactly that truncation, and the fallback ($\text{age}+1$) guarantees a
well-defined one-year life instead of a zero or negative lifespan.}

\why{\emph{A consequence that must be stated plainly.} An agent is removed from the population only at its
natural death age. A recorded cancer death does \emph{not} remove the agent; cancer deaths are counted as
cause-specific events in the statistics layer. The age structure therefore follows the all-cause life table
exactly, and the at-risk denominator used for cause-specific rates still contains agents who already have a
cancer death recorded. This is the standard device of estimating a cause-specific hazard with the
general-population denominator, and it is correct under the usual assumption that the two causes act
independently --- that removing one cause does not change the force of the other. It would be the wrong
device if a cancer death were assumed to change the subsequent risk of dying from other causes, and it is
one of the assumptions a validation exercise should probe.}

\alt{\begin{itemize}[nosep]
  \item \textbf{An annual Bernoulli draw with a running survival flag.} Distributionally equivalent for a
        single cause, but it makes the death age an output rather than an input, which makes the paired
        counterfactual much harder, and it adds a state variable that must be kept consistent.
  \item \textbf{A parametric mortality law (Gompertz--Makeham, Heligman--Pollard, Lee--Carter).} Smooth,
        compact, extrapolable; unnecessary here because the whole age range (0--110) is observed, so a
        fitted law could only add misfit.
  \item \textbf{Cause-deleted (other-cause) life tables.} Conceptually the right construction for a cancer
        model: remove cancer deaths from the background so the two causes are not double-counted. Not used
        because the packaged aggregate file supplies all-cause deaths only --- the required input is a
        cause-specific death column by age. The helper that converts annual risks into a death-age
        distribution (\id{risks\_to\_distribution}, $P(T=a)=q(a)\prod_{j<a}(1-q(j))$) is present in the
        code but never called: the model uses observed deaths directly.
  \item \textbf{Frailty models (a Gamma-distributed individual hazard multiplier).} Captures mortality
        selection --- the frail die first, so survivors are a selected subgroup --- and is a
        well-established correction. Not used because the frailty distribution is not identifiable from
        aggregate death counts, and it would quietly change the meaning of every age-specific rate.
  \item \textbf{Explicit competing-risks hazards with a copula.} The most complete treatment, allowing the
        two causes of death to be dependent; requires information on the degree of dependence that
        aggregate cross-sectional data cannot provide.
\end{itemize}}

\section{Cancer Onset: The Diagnosis Hazard}

\subsection{Functional form of the hazard}

\mech{The annual probability that a cancer is clinically diagnosed at age $a$ (given no earlier
diagnosis) is a \emph{piecewise-constant hazard on the log scale}:
\[
\log h(a) \;=\; \beta_0 + \sum_{k=1}^{K}\beta_k\,\mathbf{1}\{a \ge t_k\},
\]
with break-point ages $t_k$ and coefficients $\beta_k$ supplied by the configuration (the packaged
default is $t = (-1,0,40,50,60,70)$ with $\beta = (-11.9,0.01,0.28,-0.11,-0.1,-0.1)$). The value of
$h(a)$ is capped at 1. The \emph{diagnosis age} of an agent is then the first age $a$ at which a
Bernoulli trial with probability $h(a)$ succeeds, scanning ages from 0 upward; if no age succeeds, the
agent is never diagnosed and receives a sentinel diagnosis age ($110$, the ``unreal life length'') that
means ``no clinical cancer''. The scan is vectorised: for every agent, a column of uniform draws is
compared with the hazard curve and the first success is taken.}

\why{The functional form follows directly from what can be calibrated and what the model needs next.

\emph{(i) The calibration target is an age-specific annual incidence rate}, and a piecewise-constant
discrete hazard is exactly the quantity that corresponds to it: ``the annual probability of a diagnosis
given no diagnosis yet''. Choosing a hazard (rather than a distribution of diagnosis ages) means the
fitting step is linear in the coefficients and can be solved in milliseconds, which is what makes an
interactive web interface possible.

\emph{(ii) The break points are placed where the observed curve changes slope.} Incidence rises steeply
from mid-life, and the knots at 40, 50, 60 and 70 are the ages where the slope changes in the target
data, so a handful of parameters captures the shape without imposing a functional family.

\emph{(iii) The log scale does two jobs.} It keeps $h$ strictly positive without constraints, and it lets
the constant term express a baseline probability over several orders of magnitude --- the fitted
$\beta_0 \approx -11.9$ corresponds to a hazard of order $10^{-5}$ per year at young ages, which is the
right order of magnitude for a rare cancer and cannot be represented comfortably on a linear scale.

\emph{(iv) The cap at 1 and the first-success scan make the object a proper probability and a proper
first-event.} Capping is required because the Bernoulli device needs a probability; taking the first
success encodes the fact that ``age at diagnosis'' is a first-event time, so no agent can be diagnosed
twice and no agent can accumulate two independent diagnoses.}

\alt{\begin{itemize}[nosep]
  \item \textbf{A smooth parametric hazard (Gompertz, Weibull, log-logistic, gamma).} One or two
        parameters, a smooth derivative, easy to communicate; the risk is systematic misfit precisely
        where it matters, at the turnover and in the late-life plateau of cancer incidence. It is the
        natural choice if the target is a single narrow age range rather than the whole lifespan.
  \item \textbf{Splines (natural cubic, B-spline, or a penalised P-spline).} More flexible than step
        functions and smoother than a piecewise-constant fit; costs extra parameters and a penalty
        weight to choose. This is the most natural refinement of the present design.
  \item \textbf{Using the observed age-specific rates directly as the hazard} (no fitting at all).
        Removes the identifiability question entirely; the cost is sampling noise, no smoothing, no
        ability to interpolate or extrapolate, and a subtle mismatch between an annual rate and an annual
        discrete hazard unless the person-years correction is applied.
  \item \textbf{A mechanistic carcinogenesis model} (two-stage clonal expansion / Moolgavkar--Venzon--Knudson,
        Armitage--Doll multi-stage, or multistage with genomic instability). Explains \emph{why} incidence
        rises roughly as a power of age and links initiation and promotion to exposures, which would make
        risk-factor modelling possible. Not used because such models are not identifiable from
        age-specific incidence alone --- many parameter combinations reproduce the same incidence curve
        --- and because the extra structure does not improve the screening question being asked.
  \item \textbf{Age--period--cohort modelling.} The right tool for incidence \emph{trends}; it needs
        several periods and cohorts, whereas this model deliberately simulates a single cross-section.
  \item \textbf{Separate hazards for progressive and non-progressive lesions} --- the adenoma--carcinoma
        sequence in colorectal cancer, with a pre-clinical detectable phase of its own. This is the
        correct structural treatment of pre-cancer and is the single largest structural omission in the
        present model (Section~\ref{sec:assumptions}); it requires lesion-level data.
\end{itemize}}

\subsection{What counts as ``having cancer'', and the oldest agents}

\mech{An agent is marked \id{has\_cancer} if and only if the sampled diagnosis age is below the sentinel
($110$) \emph{and} below the agent's natural death age. A person who would die before their clinical
diagnosis year is therefore never given a tumour history at all. In addition, agents who are already
older than 98 at $t=0$ are assigned the sentinel directly, so they can never develop cancer during the
run.}

\why{Both rules follow from the same decision: the model's unit of cancer is \emph{a cancer that would be
clinically diagnosed in this person's remaining lifetime}. That is exactly the event that produces the
observed incidence and mortality, and exactly the event that screening could displace. An agent who dies
before the diagnosis year contributes nothing to any reported quantity, so constructing a tumour history
for them would add cost and a meaningless trajectory. The age guard has the same motive in a cruder form:
the diagnosis-age scan is bounded by the maximum representable age, and back-calculating a preclinical
onset age for an agent already past 98 would place the entire preclinical trajectory outside the
simulated window, so the agents concerned are simply excluded.}

\why{\emph{What this costs, stated honestly.} The natural history is \emph{diagnosis-driven}: everybody
who has cancer in the model would also have been diagnosed clinically. The model therefore cannot
represent a cancer that exists but never becomes clinically apparent --- which is precisely the
overdiagnosis mechanism. Without screening this is harmless, because such a cancer is invisible in every
reported outcome anyway. Under a screening scenario it is not harmless: every screen-detected cancer is a
``true'' cancer that would have been diagnosed and could have killed, so in the presence of a real
non-progressive fraction the model will overstate the benefit of screening. The age guard is also a blunt
modelling shortcut with no epidemiological justification --- it is a range guard dressed as a risk rule
--- and it should be revisited if the horizon, the initial age structure or the sentinel age changes.}

\alt{\begin{itemize}[nosep]
  \item \textbf{Onset-driven natural history.} Give every agent a preclinical onset (from an initiation
        hazard) and let progression, symptoms and screening determine whether and when it is ever
        diagnosed. Then a person can die with undiagnosed cancer, and both overdiagnosis and
        ``screen-detected but would never have surfaced'' become representable. This is the structurally
        correct design and the direction a serious upgrade should take; it requires a progression model
        and preclinical-detectability data that the packaged inputs do not contain.
  \item \textbf{A non-progressive (or slowly progressive) fraction.} Standard in screening models: a
        fraction $\mathrm{nc}$ of lesions never reach clinical significance. One extra parameter, and it
        is the accepted way to make overdiagnosis quantifiable --- but it cannot be estimated from
        incidence and mortality data alone, so it would have to be imposed from the literature and made a
        target of sensitivity analysis.
  \item \textbf{No exclusion at all: build a tumour history for every agent regardless of its fate.}
        Conceptually uniform and it removes the special case, but it multiplies the work for information
        that no output consumes, and it produces ``cancer in people who are already dead''.
  \item \textbf{Removing the 98-age guard but widening the sentinel.} A pure implementation fix: the
        guard should either be dropped (letting the hazard decide) or replaced by an explicit, documented
        truncation age.
\end{itemize}}

\section{The Preclinical Phase and the Time-Scale Problem}

\subsection{Exponential lead time}

\mech{For every agent with cancer the model draws a \emph{lead time} (preclinical sojourn time, i.e.\ the
time from the first malignant cell to clinical diagnosis) from an exponential distribution with mean
\id{lead\_time\_mean} (3.5 years in the default configuration), so
$t_{\text{ttd}} \sim \mathrm{Exp}(\lambda)$, $\lambda = 1/\id{lead\_time\_mean}$. The onset age is then
back-calculated from the \emph{observed} diagnosis age: $\text{onset} = \max(\text{diag} - \lceil t_{\text{ttd}}\rceil, 0)$.
The continuous value of $t_{\text{ttd}}$ is used for the stage determination (Stage~\ref{sec:stage}) and the
integer value for the age arithmetic.}

\why{The whole preclinical history is anchored to the diagnosis date rather than simulated forward from an
onset date, and that choice is the reason the model can be calibrated from aggregate data at all.

\emph{(i) The observable is the diagnosis age.} Anchoring the preclinical trajectory to it means the model
reproduces the observed incidence exactly by construction and uses the lead time only to decide what the
tumour looked like \emph{before} a date that is already correct. Nothing has to be right about the absolute
onset rate, so a whole class of misfit is removed.

\emph{(ii) Exponential is the maximum-entropy choice with a given mean.} Assuming a constant progression
rate --- no memory of how long the tumour has already been present --- is the least committal assumption
available for a sojourn time, and it is exactly what the classical Markov screening literature assumes.
That compatibility is not cosmetic: the relations between lead time, sojourn time and screening benefit
derived in that literature remain usable to sanity-check the model's behaviour.

\emph{(iii) One parameter, because only one can be afforded.} The absolute scale of the preclinical phase
is not identifiable (Section~\ref{sec:nonident}), so adding a second sojourn parameter would add a second
unidentified quantity. The mean is the smallest sufficient summary: it is the scale of everything
downstream (crossing times, incidence age, stage duration) and it is the single knob the sensitivity
analysis turns.

\emph{(iv) Memorylessness keeps the staging calibration closed-form.} Under an exponential lead time
$P(\text{stage} \ge s) = \exp(-t_s/\mu)$, which is what makes it possible to \emph{invert} the observed
stage mix into crossing times with one logarithm per stage (Section~\ref{sec:stage}) instead of solving a
numerical calibration problem.

\emph{(v) Continuous for the comparison, integer for the arithmetic.} Rounding the lead time before
comparing it with the crossing times would bias the stage distribution (a rounding-down would
systematically over-produce early stages), so the crossing comparison uses the continuous draw while the
age arithmetic uses its ceiling.}

\alt{\begin{itemize}[nosep]
  \item \textbf{Gamma or Weibull sojourn times.} The standard generalisations, widely used in
        screen-detection models; exponential is their shape-$1$ special case. They buy a more realistic
        shape (few very long sojourns), but the shape parameter has to be assumed rather than estimated,
        so they trade one assumption for two.
  \item \textbf{A deterministic (fixed) lead time.} The simplest possible assumption, and the stage mix
        would then be a knife-edge function of a single threshold rather than a smooth probability; it
        also makes the preclinical duration unrealistically discrete.
  \item \textbf{Per-stage dwell times} (a sojourn time for each stage). Much closer to tumour biology and
        it makes stage duration directly interpretable; the available data cannot inform four separate
        dwell times, and it forces the crossing times onto an integer grid.
  \item \textbf{Size-driven detectability from the growth curve itself.} If the tumour is ``detectable''
        once it crosses a size threshold, the sojourn time is whatever the Gompertz model says it is,
        which is beautifully self-consistent because the model already contains that curve. The missing
        inputs are a detectable-size threshold and a size-dependent test sensitivity, neither of which is
        in the packaged data.
  \item \textbf{Estimating the lead time from screening data} (screen-detected versus interval cancers,
        or a randomised screening trial). This is the only route to genuine identification, and it is the
        single highest-value data addition the model could receive.
  \item \textbf{Anchoring the mean to the literature} instead of leaving 3.5 as a free number, with the
        published range used as the sensitivity interval. Operationally identical to the present design,
        but it turns an assumption into a citable assumption --- the cheapest available improvement.
\end{itemize}}

\subsection{Non-identifiability of the time scale, and what is fixed instead}
\label{sec:nonident}

\mech{The model does \emph{not} estimate the preclinical time scale. It fixes the lead-time mean
(\id{lead\_time\_mean}) and calibrates everything else and, separately, derives the stage-crossing times
from the observed stage mix (Section~\ref{sec:stage}). The fixed scale is then exposed as the single axis
of the sensitivity analysis (Section~\ref{sec:sens}).}

\why{The reason is a genuine mathematical property of the system, not a shortage of effort. Let
$c_s = C - \ln(K - \ln s)$ and write the crossing time of stage $s$ as $t_s = c_s / B_{\text{ind}}$
(Section~\ref{sec:gompertz}). Under the exponential lead time with mean $\mu$,
\[
P(\text{diagnosed at stage} \ge s) \;=\; P(t_s \le t_{\text{ttd}}) \;=\; \exp\!\left(-\frac{c_s}{B_{\text{ind}}\,\mu}\right),
\]
so the observed stage mix constrains only the \emph{product} $B_{\text{ind}}\mu$. Doubling the growth rate
and halving the sojourn mean leaves every observable exactly as it was. Formally, the parameter vector
$(B_{\text{ind}}, \mu)$ is unidentified along the hyperbola $B_{\text{ind}}\mu = \text{const}$ --- a flat
likelihood ridge, and a likelihood maximiser would happily return an arbitrary point on it.

There is a second, more basic reason. Onset --- the appearance of the first malignant cell --- is not an
observable event. The lead time is therefore not a quantity that is measured badly; it is a quantity that
is, in principle, unobservable in a clinical dataset, because diagnosis interrupts the very history one
would want to observe. ``Estimated lead time'' in the literature always means ``estimated under an
explicit model''.

What the data \emph{do} identify is: the diagnosis hazard (from age-specific incidence), the product
$B_{\text{ind}}\mu$ (from the stage mix), and the terminal-phase hazard (from cancer mortality). The design
choice is therefore to spend the one remaining degree of freedom as an \emph{explicit assumption with an
interpretable name} rather than as a fitted value that pretends to be an estimate. Fixing $\mu$ rather than
$B$ is the asymmetric but correct choice: $\mu$ has literature estimates for colorectal cancer and a plain
clinical meaning (time from onset to diagnosis), whereas $B$ has a scale that nothing outside the model
can speak to.

Concretely, the consequence is a clean division of labour: $\mu$ is an assumption, the crossing times are
\emph{derived} from the data, and the growth-rate parameters are \emph{fitted} to those derived times. So
the model reproduces the observed stage mix by construction, and the single assumption is visible in the
parameter file, in the document, and in the sensitivity plots.}

\alt{\begin{itemize}[nosep]
  \item \textbf{Report the lead-time range as the headline result, not as a sensitivity appendix.} Nothing
        about the maths changes; what changes is that the reader sees a family of curves as the answer
        rather than one curve plus a caveat. For a screening question this is arguably the more honest
        presentation, and it is a reporting change rather than a modelling change.
  \item \textbf{Bayesian treatment: a prior on $\mu$, the rest fitted, and a full posterior.} The
        statistically cleanest way to say ``$\mu$ is a judgement'' --- the judgement enters as a prior,
        and its consequences arrive as credible intervals on every output rather than as a set of
        re-runs. Requires MCMC (or ABC) at a cost of hours instead of milliseconds, and it makes the
        headline numbers depend on a prior that is itself not derived from data.
  \item \textbf{Profile likelihood over $\mu$.} Compute the fit as a function of $\mu$ and show the
        surface. Cheap (the model is fast), transparent, and it makes the ridge visible rather than
        merely described. This is the smallest possible step towards a proper uncertainty treatment.
  \item \textbf{Fix $B$ instead of $\mu$.} Symmetric and equally arbitrary, and strictly worse: the growth
        rate has no external anchor at all, so the choice would be uninterpretable and untestable.
  \item \textbf{Two-parameter sojourn time.} Worse, not better: it replaces one unidentified scale with
        two unidentified parameters, and no amount of prior opinion can identify them.
  \item \textbf{External identification from screening data.} A screening trial or a registry of
        screen-detected and interval cancers identifies the sojourn time directly; that is the only way to
        turn $\mu$ into an estimate. It requires individual-level screening data this model does not have.
\end{itemize}}

\section{Tumour Growth: The Reduced Gompertz Model}
\label{sec:gompertz}

\subsection{The growth curve}

\mech{Tumour growth is described by the \emph{reduced Gompertz} function
\[
V(t) = \exp\!\bigl(K - \exp(C - B\,t)\bigr),
\]
where $V$ is an abstract proxy for tumour mass (a ``size''), $t$ is years since the first malignant cell,
$K$ and $C$ are shape parameters and $B$ is the growth rate. The curve starts at
$V(0) = \exp(K - \exp C)$ (a small tumour), saturates at $V(\infty) = \exp K$ (the carrying capacity), and
has its inflection point at $t = C/B$, where the growth switches from accelerating to decelerating. Stages
are the crossing times of the sizes $2$, $3$ and $4$:
\[
t_s = \frac{C - \ln(K - \ln s)}{B}, \qquad s \in \{2,3,4\},
\]
which is finite exactly when $\exp K > s$, i.e.\ when the carrying capacity is above the stage threshold.
In the reference fit ($K = 1.451$, $C = 0.198$, $B = 0.550$, so $\exp K = 4.27$) the crossing times are
$0.87$, $2.26$ and $5.34$ years.}

\why{Gompertz is chosen for four reasons, three substantive and one structural.

\emph{(i) It is the classical tumour growth law and it matches the qualitative shape of the data.} Tumours
grow roughly exponentially at first and then decelerate as they outstrip their blood supply and available
space; the Gompertz curve has exactly that feature --- exponential in the early phase, asymptotically flat
late --- and has been the default descriptive growth model in oncology for half a century.

\emph{(ii) It is parsimonious without being rigid.} Two shape parameters ($K$, $C$) plus one rate ($B$)
reproduce a wide range of growth curves, and each parameter has a plain meaning: $\exp K$ is a carrying
capacity in size units, $C$ fixes the initial size through $V(0) = \exp(K - \exp C)$, and $B$ is the
specific growth rate.

\emph{(iii) Its crossing times have a closed form.} The inverse
$t_s = (C - \ln(K - \ln s))/B$ is what makes the staging calibration a \emph{linear-in-a-logarithm}
problem and lets the model store each agent's stage ages at generation time. A growth law whose stage
times required numerical inversion would be both slower and much harder to reason about --- and the
closed form is also what makes the non-identifiability statement of Section~\ref{sec:nonident} provable
rather than empirical.

\emph{(iv) Crossings of absolute sizes $\{2,3,4\}$ are a normalisation, not a restriction.} The model has
three crossing times as calibration targets and three free parameters ($K$, $C$, $B$), so the fit is
exactly identified: the threshold convention is absorbed by the fitted shapes. What the threshold spacing
\emph{does} affect is the individual-level dispersion of stages (through $B_{\text{std}}$,
Section~\ref{sec:hetero}), because different tumours cross differently spaced thresholds at different
times. That is a limitation to keep in view, not a hidden distortion of the population means.}

\alt{\begin{itemize}[nosep]
  \item \textbf{Exponential growth.} Simplest possible law; unbounded, so it has no carrying capacity, and
        it forces the crossing times to be equally spaced in log size --- a strong restriction that the
        data would contradict.
  \item \textbf{Logistic growth.} Familiar and symmetric; the deceleration is symmetric about the
        inflection, which is known to be wrong for tumours (they grow for a long time before slowing),
        and its crossing times are symmetric for the same reason.
  \item \textbf{von Bertalanffy / Bertalanffy--Pütter.} Mechanistically motivated (growth as the balance of
        anabolic and catabolic terms) and in the same family as Gompertz; comparable behaviour, less
        standard in the screening literature.
  \item \textbf{Weibull or power-law growth.} Flexible; no carrying capacity, and no established clinical
        interpretation of the parameters.
  \item \textbf{A mechanistic, nutrient- or space-limited ODE model with explicit vasculature.} The most
        biologically grounded option; many parameters, and it needs volumetric growth data that aggregate
        cancer statistics do not contain.
  \item \textbf{Stochastic branching-process growth.} Models the clonal dynamics directly and --- uniquely
        --- can express \emph{extinction} of early tumours, which is a genuine mechanism of
        non-progression. Computationally far heavier and not identifiable from the available data.
  \item \textbf{Piecewise-linear progression (no growth model).} A constant progression rate per tumour
        class, i.e.\ the classic multi-state Markov formulation. Simpler and directly fittable to stage
        data; it severs the link between growth and stage timing that makes the current model
        interpretable, and it leaves the notion of a ``size'' with nothing behind it.
  \item \textbf{Direct ordinal regression of stage on covariates} --- the design this replaced. It
        reproduces the observed stage mix, but the stage becomes an arbitrary draw decoupled from any
        growth process, individual heterogeneity loses its physical meaning, and the failure is silent
        because the outputs still look plausible.
\end{itemize}}

\subsection{Individual heterogeneity in the growth rate}
\label{sec:hetero}

\mech{Each agent with cancer receives its own growth rate
\[
B_{\text{ind}} = \exp\!\bigl(\ln B_{\text{pop}} + \varepsilon\bigr), \qquad
\varepsilon \sim \mathcal{N}(0, B_{\text{std}}^2),
\]
so $B_{\text{ind}}$ is log-normally distributed with median $B_{\text{pop}}$ and multiplicative spread
$\exp(B_{\text{std}})$ (a floor of $10^{-6}$ prevents a degenerate zero). The population-mean curve used
for fitting uses $B_{\text{pop}}$ without the random effect. $B_{\text{std}}$ is \emph{not} fitted: it is a
hand-set value (0.1 by default, i.e.\ roughly $\pm10\%$ multiplicative spread).}

\why{Two arguments, one about realism and one about a specific screening phenomenon.

\emph{(i) Growth rates really do vary, and the clinical fact is not a nuisance parameter.} Tumour
doubling times for the same cancer differ between patients by factors, not percentages. Without
heterogeneity, the stage is a deterministic function of the lead time: every patient with the same
elapsed preclinical time would sit at the same stage, which imposes a rigid, artificial correlation
between ``how long the tumour has been growing'' and ``how advanced it is''. Heterogeneity decouples them,
so two patients with identical lead times can legitimately be diagnosed at different stages.

\emph{(ii) Length bias needs it.} Screening preferentially catches slowly growing tumours, because a
tumour that spends longer in the detectable preclinical phase has more chances to be caught. This
\emph{length bias} --- together with the associated lead-time bias --- is one of the classical reasons
screening trials overstate benefit, and it can only appear in a model in which growth rates differ between
individuals. With $B_{\text{std}} = 0$ the phenomenon vanishes by construction and the model's stage shift
becomes too favourable and too clean.

Log-normal rather than normal is the natural parameterisation for a rate: positivity is automatic, the
distribution is right-skewed (a few fast tumours, a floor of slow ones) as multiplicative quantities tend
to be, and the two parameters have direct meaning --- $B_{\text{pop}}$ is the median and $B_{\text{std}}$
the log-scale spread. One subtlety worth remembering: the arithmetic mean is
$B_{\text{pop}}\exp(B_{\text{std}}^2/2)$, slightly above the median, so a mean-based comparison against a
single fitted rate carries a small, known offset.}

\why{\emph{An honest weakness.} $B_{\text{std}}$ governs how much individual variation, stage dispersion
and length bias the model produces --- and it is set by hand, not estimated. Nothing in the current
calibration constrains it, because the fitting step targets only the population-mean crossing times
($K$, $C$, $B_{\text{pop}}$). It is therefore a second unverified assumption alongside the lead-time mean,
but unlike the lead time it is currently \emph{not} varied in the sensitivity analysis. Turning it into a
sensitivity axis is cheap and would materially strengthen the uncertainty statement.}

\alt{\begin{itemize}[nosep]
  \item \textbf{No heterogeneity ($B_{\text{std}} = 0$).} The simplest option and the one that makes the
        stage a deterministic function of the lead time. It removes length bias and stage dispersion
        entirely and makes the correlation between stage, lead time and survival artificially tight ---
        so it should only be used as a deliberate sensitivity comparison, never as the baseline.
  \item \textbf{A Gamma-distributed growth rate.} Also positive, with independent control of shape and
        scale; less natural for a multiplicative quantity and no better suited to the data.
  \item \textbf{A discrete number of growth classes} (slow / medium / fast). Very easy to communicate and
        to vary in a scenario analysis; coarse, introduces artificial discreteness, and the class
        boundaries are arbitrary.
  \item \textbf{A normal (untruncated) growth rate.} Rejected: it admits negative rates, which are not
        merely unlikely but meaningless.
  \item \textbf{A full random-effects model of $\ln B$ with covariates} (age, sex, stage at onset). The
        right way to make growth heterogeneity explicit rather than incidental; needs covariate-specific
        growth information that the staging file does not contain.
  \item \textbf{Empirical volume-doubling times} from imaging series or from the literature. The only
        route to a defensible value of $B_{\text{std}}$; it would convert a hand-set number into a cited
        one.
\end{itemize}}

\subsection{From growth to stage: staging as a derived quantity}
\label{sec:stage}

\mech{For each agent with cancer the model computes its own crossing times $t_2, t_3, t_4$ (from
$B_{\text{ind}}$, Section~\ref{sec:gompertz}) and compares them with the continuous lead time: the stage
at clinical diagnosis is $4$ if $t_4 \le t_{\text{ttd}}$, else $3$ if $t_3 \le t_{\text{ttd}}$, else $2$
if $t_2 \le t_{\text{ttd}}$, else $1$. In other words the stage is \emph{the number of size thresholds
the tumour has crossed by the time it is clinically diagnosed}. The full stage-age array stored per agent
is $[\text{onset},\ \text{onset}+t_2,\ \text{onset}+t_3,\ \text{onset}+t_4,\ \text{onset}+t_4+\delta]$,
where $\delta$ is an exponential draw for the terminal phase after the last crossing.

The crossing times used as calibration \emph{targets} are obtained by inverting the observed stage mix.
Because the lead time is exponential, $P(\text{stage} \ge s) = \exp(-t_s/\mu)$, so
\[
t_2 = -\mu \ln(1 - p_1), \qquad t_3 = -\mu \ln(p_3 + p_4), \qquad t_4 = -\mu \ln p_4 ,
\]
where $p_s$ is the observed proportion at stage $s$ and $\mu$ is the lead-time mean. For the USA
colorectal staging file the observed mix is $0.220/0.256/0.308/0.218$, which yields the targets
$0.87/2.26/5.34$ years; the resulting fit reproduces the mix with $0.225/0.250/0.284/0.241$
($N = 20\,000$, seed 2022).}

\why{Staging is derived rather than drawn, and the difference matters more than it looks.

\emph{(i) It makes the stage mix a consequence instead of an input.} When the stage is drawn from a
distribution, reproducing the observed mix is a tautology --- the model was told the answer. Here the mix
is produced by the interaction of two physical quantities (how long the tumour has been present, how fast
it grows), so agreement with the observed mix is a genuine, falsifiable calibration check. The residual
gap ($0.225/0.250/0.284/0.241$ versus $0.220/0.256/0.308/0.218$) is small and attributable to sampling
noise, to $B_{\text{std}}$, and to the interaction between the stage and the age-dependent diagnosis
hazard --- not to a systematic misfit.

\emph{(ii) It is what makes the per-agent growth rate meaningful.} If the stage came from a regression,
$B_{\text{ind}}$ would be a number with no downstream consequence; because the crossing times are computed
from it, the growth rate controls the stage, the stage ages, and therefore the timing of the terminal
phase. Growth heterogeneity has observable consequences, which is exactly what a structural parameter
should have.

\emph{(iii) The inversion is closed-form, so the calibration target is data-derived and auditable.} One
logarithm per stage turns the staging file into three numbers that can be printed, checked by hand, and
traced back to the source table. There is no inner optimisation loop to hide a mistake in.

\emph{(iv) It respects the non-identifiability accounting.} The inversion consumes $\mu$ (the assumed
scale) and returns the crossing times; the Gompertz fit then consumes the crossing times and returns
$K$, $C$, $B_{\text{pop}}$. Three targets, three parameters --- an exactly identified chain in which the
single assumption is visible at the front.}

\alt{\begin{itemize}[nosep]
  \item \textbf{Draw the stage directly from the observed distribution.} Reproduces the mix exactly, at the
        cost of making the stage exogenous: growth heterogeneity, lead time and screening can no longer
        shift the stage except by an explicit, arbitrary fudge, so the tumour growth model becomes
        decorative.
  \item \textbf{Multinomial or ordinal regression of stage on age, sex and other covariates.} The design
        this replaced. It fits the mix well and is easy to implement, but it decouples the stage from the
        growth process, leaves the individual growth rate meaningless, and its coefficient interpretation
        (``the effect of age on the log-odds of being at a later stage'') has no mechanism behind it.
  \item \textbf{A latent stage-progression process with transition rates} (the standard multi-state
        formulation). Clean, and it makes dwell times explicit; it needs observed stage-transition data
        (e.g.\ from serial screening rounds), which the packaged inputs lack.
  \item \textbf{Clinical TNM staging} (tumour size $T$, nodal involvement $N$, distant metastasis $M$).
        The clinically correct and far more informative description, and it would allow a genuine
        screening stage shift to be measured against real thresholds; it requires the size and nodal
        information that a stage-only staging file does not contain.
  \item \textbf{Discrete growth classes with fixed dwell times} (slow/intermediate/fast tumours,
        each with a fixed dwell time per stage). Simple and transparent, and it is how several published
        screening models are parameterised; it makes the dispersion of stages a discrete choice rather
        than a continuous one.
\end{itemize}}

\section{Aggressiveness: A Cure-Modifying Tumour Label}
\label{sec:agg}

\subsection{What the label is and what it does}

\mech{Every cancer agent is labelled \emph{aggressive} or not by a Bernoulli draw whose probability is
taken from the staging data: for a tumour diagnosed at stage $s$ the probability is the observed
proportion of aggressive tumours in that stage, computed from the file's \id{Aggressiveness} column
(USA colorectal data: $0.096 / 0.198 / 0.437 / 0.639$ for stages I--IV, i.e.\ aggressiveness and stage
advance together). The label then affects \emph{exactly one thing} --- the probability of cure. The
baseline cure probability $p$ is converted to odds, multiplied by a per-stage odds ratio
(\id{aggressiveness\_cure\_odds\_ratio}, default $0.667$ for every stage), and converted back:
\[
\text{odds}' = \frac{p}{1-p}\times \mathrm{OR}_s, \qquad
p' = \frac{\text{odds}'}{1 + \text{odds}'}.
\]
A value of $\mathrm{OR}_s = 1$ is a strict no-op (the transform is skipped), so the assumption can be
switched off without touching any other behaviour.}

\why{Four reasons, and they interlock.

\emph{(i) It is the only phenotype dimension for which the model has data.} The staging file carries an
aggressiveness flag, so the per-stage prevalence is \emph{measured}, not assumed. Ignoring it would throw
away real information; inventing a different tumour-phenotype axis would have no data behind it at all.

\emph{(ii) It is deliberately not a growth modifier, to avoid double counting.} Aggressive tumours are, in
clinical intuition, fast-growing ones --- so the tempting design is to let the label change $B$. That would
confound it with $B_{\text{std}}$, which already carries growth heterogeneity, and the calibrated stage mix
would then have to absorb two overlapping sources of variation. Assigning the label to the cure
probability puts it where it is observable: in survival, not in staging.

\emph{(iii) The odds-ratio form is the right algebra for a cure probability.} Cure is a binary outcome, so
an odds ratio is the standard effect measure; it composes multiplicatively with the baseline odds, cannot
leave $[0,1]$, is symmetric in framing ($\mathrm{OR}$ and $1/\mathrm{OR}$ express ``harder'' and
``easier'' with the same parameter), and reduces to the identity at $\mathrm{OR} = 1$.

\emph{(iv) The direction is a clinical prior, and the form leaves room for stage-specific evidence.} That
an aggressive tumour is harder to cure is a reasonable prior, and $\mathrm{OR}<1$ encodes it. The parameter
is a vector over stages even though the default sets all four equal, so a user with stage-specific
treatment-response evidence can enter it without changing the design.}

\why{\emph{Two limitations to state plainly.}
First, the odds ratio itself is \emph{not estimated from data} --- it is set by hand. The staging file
contains no treatment-outcome column, so there is nothing in the packaged inputs that could identify it:
what is measurable (the prevalence by stage) is measured, and what is not (the effect of the label on
cure) is an explicit assumption with a clinical prior.
Second, because aggressiveness is strongly associated with stage ($0.096$ versus $0.639$), a pooled
comparison of aggressive against non-aggressive survival confounds the odds ratio with the stage mix ---
the aggregate numbers will show a large survival gap that is mostly a stage effect. This is a
Simpson's-paradox trap, and the model avoids the trap by reporting survival curves and cure fractions
\emph{stratified by stage and aggressiveness} rather than only pooled (Section~\ref{sec:survival}).}

\alt{\begin{itemize}[nosep]
  \item \textbf{Aggressiveness affecting growth ($B$).} Biologically natural, and it would make the label
        do something visible in the staging data. Rejected because it double-counts against
        $B_{\text{std}}$ and destroys the exact stage-mix calibration.
  \item \textbf{A continuous latent aggressiveness score.} More realistic --- aggressiveness is a
        spectrum, not a flag --- and it would let the cure odds vary smoothly; it cannot be derived from a
        binary column, and it would add a non-identifiable latent dimension.
  \item \textbf{A mixture cure model with the label as a covariate.} Statistically the standard
        formulation, and the present design is effectively its logistic special case; a full mixture model
        would estimate the uncured fraction directly, at the cost of a likelihood that needs treatment
        follow-up data.
  \item \textbf{Subtype-specific (grade, MSI/MMR status, molecular subtype) survival from the literature.}
        Substantially more informative and clinically meaningful than a binary flag, and it would make the
        cure effect evidence-based. Requires subtype-level data the aggregate files do not contain.
  \item \textbf{Making aggressiveness modify the cancer-death hazard instead of the cure probability.}
        Equivalent in its effect on mortality, different in mechanism (it changes survival given
        non-cure rather than the chance of cure). Arguably more direct; harder to calibrate and less
        interpretable as ``treatment success''.
  \item \textbf{Dropping the label.} Simplest, and it removes a hand-set parameter --- but it discards the
        only tumour-phenotype information the inputs provide.
\end{itemize}}

\section{Treatment and Cure}

\subsection{The cure probability}

\mech{At clinical diagnosis the model computes a cure probability
\[
p_{\text{cure}} \;=\; \min\!\Bigl(1,\; \alpha(\text{diag})\times \theta_{s}\Bigr),
\]
where $\theta_s$ is the stage-specific treatment efficiency and $\alpha(\cdot)$ is an age-dependent
efficiency factor. $\theta = [0.7,\,0.3,\,0.2,\,0.2]$ for stages I--IV. $\alpha$ is a step function of the
age at diagnosis with break points $40/50/60/70$ and values $[1.0,\,0.9,\,0.5,\,0.3,\,0.1]$. If the tumour
is aggressive, the cure \emph{odds} are multiplied by the per-stage odds ratio of
Section~\ref{sec:agg} before converting back. Cure is then a Bernoulli draw with the resulting
probability. In the reference configuration $p_{\text{cure}}$ spans $0.70$ (stage I, diagnosed young) down
to $0.02$ (stage IV, diagnosed after 70).}

\why{The cure probability is factored into stage and age because those are the two axes on which
curability is actually observed and reported.

\emph{(i) Stage dominates and the data measure it.} Stage at diagnosis is the strongest single predictor of
cancer survival in every registry, and it is the dimension along which the staging file is organised.
Making it the primary factor aligns the model with the observable that constrains it.

\emph{(ii) Age enters as a step function deliberately.} The age gradient in cure is real --- older patients
tolerate treatment worse, are treated less aggressively and carry more competing mortality --- but a smooth
age effect would need shape parameters that nothing in the inputs can identify. A four-step function is
the granularity at which cancer statistics are published and at which the outputs are reported, it needs
no shape assumption, and it is directly legible to a clinician. The step values read as ``relative cure
efficiency in this age band''.

\emph{(iii) Multiplication is the right composition.} It keeps the result inside $[0,1]$ without clamping
except at the top, it makes each factor interpretable as a relative multiplier (a patient with the best
possible age efficiency receives exactly the stage-specific efficiency), and it makes switching an effect
off a one-line change (all age constants set to $1$).

\emph{(iv) The cap at $1$ is a correctness requirement} --- the Bernoulli device needs a probability --- and
it is harmless because the reference values never approach it.

\emph{(v) The odds-ratio transform for the aggressive label behaves sensibly.} It is monotone and bounded,
so the label's effect is largest at intermediate cure probabilities and shrinks towards zero at both
extremes, where a cure probability is already nearly certain or nearly hopeless.}

\subsection{What ``cured'' means for death}

\mech{A cured patient receives a cured death age, but the decision of \emph{when} the person dies is made
by the competing-risk rule of Section~\ref{sec:death}: the agent leaves the population at its natural
death age, and the cause is attributed in the statistics layer. Each year the model therefore compares two
stored ages and counts a cancer death when the agent has cancer, is not cured (neither clinically nor by
screening), the untreated cancer-death age equals the current age, and that age does not exceed the
natural death age. The uncured patient's cancer-death age is the last stage crossing plus an exponential
draw for the terminal phase, $\delta \sim \mathrm{Exp}(\lambda_{\text{terminal}})$, where
$\lambda_{\text{terminal}}$ is fitted to the observed cancer mortality ($0.35$ per year in the reference
fit, i.e.\ a mean terminal survival of about $2.9$ years).}

\why{Defining cure as ``the cancer no longer determines either the date or the cause of death'' is what
keeps the screening arithmetic exact and honest.

\emph{(i) The paired counterfactual stays within one agent.} Because the untreated cancer-death age is
stored for every patient, a cure is simply a change of which cause the statistics count. The lives-saved
and years-saved outputs are then genuinely within-individual differences, not differences between two
populations with different sampling noise --- which is what makes them meaningful at the sample sizes
used.

\emph{(ii) No post-cure survival model is needed.} A cured patient's remaining lifetime is the
population's remaining lifetime, which is precisely the competing-risk assumption already made for
everyone else. Introducing anything else would require a cure-specific mortality (treatment toxicity,
late relapse) and therefore data.

\emph{(iii) The terminal phase carries the timing of cancer death, and it is the one treatment-adjacent
quantity that is actually fitted.} Cancer mortality in the data is dominated by how long uncured patients
survive after diagnosis, so a single fitted hazard for the terminal phase is the most economical way to
match that observable.}

\why{\emph{The confound, restated because it matters.} Because the cure probabilities are hand-set and
only the terminal hazard is fitted, the observed cancer mortality constrains the \emph{product} of the two
rather than each separately. A scenario with more cures and faster cancer death can reproduce the same
mortality curve as one with fewer cures and slower death. Consequences: cure fractions reported by the
model should be read as scenario assumptions, not as estimates; the fitted terminal hazard absorbs
whatever the assumed cure probabilities get wrong; and only stage-specific relative survival with
follow-up could separate them. This is the second flat ridge in the model, and it deserves the same
prominence as the lead-time ridge of Section~\ref{sec:nonident}.}

\alt{\begin{itemize}[nosep]
  \item \textbf{Stage-specific relative-survival tables used directly} (the classical approach in cancer
        cost-effectiveness models). The observable becomes the input, so no cure mechanism is needed and
        published survival is reproduced exactly; the mechanism is lost, ``why'' a patient survives longer
        cannot be asked, and treatment advances can only be represented by swapping the table.
  \item \textbf{A mixture cure model with the cure fraction estimated from follow-up data.} Statistically
        the standard formulation, and the current design is essentially its logistic core; it requires
        individual follow-up and would identify the cure fractions that are currently assumed.
  \item \textbf{Excess-mortality (relative survival) hazard modelling.} Model the cancer-attributable
        excess hazard as a function of stage, age and time since diagnosis instead of a probability of
        cure. More faithful to how registry data behave, and it avoids the binary ``cured / not cured''
        abstraction; it needs population life tables by age, sex and period, and it makes the cured flag
        --- and hence the lives-saved accounting --- only approximate.
  \item \textbf{Treatment-pathway modelling} (surgery, chemotherapy, radiotherapy, each with its own
        outcome distribution). Clinically the most informative, and the natural home of several currently
        inactive parameters (complications, treatment mortality); it needs treatment data that the
        packaged inputs lack.
  \item \textbf{Post-cure mortality adjustment} (a multiplier on the natural-death hazard for cured
        patients, representing treatment toxicity or late relapse). A small, honest extension; it needs
        data to set the multiplier and would break the ``cured patients are ordinary people'' simplicity.
  \item \textbf{A strict cure that deletes the cancer death hazard entirely} and re-draws other-cause
        mortality from a cause-deleted life table. Internally the purest formulation; it requires the
        cause-deleted life table discussed in Section 3.
\end{itemize}}

\section{Screening: Schedule, Participation, Test Performance}

\subsection{Invitation schedule and participation}

\mech{A screening round takes place in simulated year $y$ if $y \ge \id{screening\_date}$ and
$(y - \id{screening\_date})$ is a multiple of \id{frequency} (defaults: the programme starts in simulated
year 10 and runs every 2 years). In a round, an agent is \emph{eligible} if it is alive, its age lies in
$[\id{start\_age}, \id{finish\_age}]$ (50--60 by default), it has not already been screen-detected, and it
has not yet been clinically diagnosed (current age below its diagnosis age). Every eligible agent takes the
test with probability \id{participation\_rate} (0.5), drawn independently in each round.}

\why{The schedule is expressed as an age window plus a period because that is how screening programmes are
actually specified and evaluated, and it is representable exactly on the annual time grid
(Section 2.2). Two eligibility rules carry real modelling content. Excluding already-detected agents keeps
the history of an agent's screening path simple and reflects the clinical reality that a person found to
have cancer leaves the screening pathway. Excluding agents already past their clinical diagnosis age
guarantees that a screen-detected cancer really is a cancer caught \emph{earlier} --- which is the only
mechanism by which screening can help in this model --- and it prevents a person from being counted both as
screen-detected and as clinically diagnosed.}

\why{Independent per-round participation is the maximum-entropy assumption given a single aggregate
attendance rate, and it is the only one that an aggregate rate can support. It also reproduces the
realistic phenomenon of intermittent attendance: a person may attend at 52 and skip 54. What it
\emph{cannot} represent is the well-documented subgroup of never-attenders (and the correlated
over-attenders). That is a genuine simplification: real programmes' benefit depends on reaching the
under-screened, and a model with independent attendance will overstate the population-level effect of a
given average participation rate.}

\subsection{Test performance}

\mech{The test is described by a true-positive rate and a false-positive rate, chosen as a pair from a
menu indexed by \id{selected\_test}: $\{50,60,70,80\}$ map to
$\text{TP} = \{0.6,0.7,0.9,0.95\}$ and $\text{FP} = \{0.05,0.07,0.1,0.15\}$. Only the true-positive rate
is actually applied: an agent with preclinical cancer tests positive with probability TP, and an agent
without preclinical cancer tests positive with probability $0.1\times\text{TP}$.}

\why{A menu of (\mbox{TP}, \mbox{FP}) operating points makes the accuracy trade-off explicit and prevents
the physically impossible combination of a perfect test: the user can only pick from the efficient
frontier that the menu offers. Expressing the test as TP/FP applied to the two subgroups (with versus
without preclinical cancer) is the classical 2$\times$2 formulation and needs no distributional assumption
about test values --- which matters because no test-value data are shipped.}

\why{\emph{Two documented inconsistencies in this area, both of which should be fixed rather than
explained.} \emph{(i)} The configured false-positive rate is inert: the false-positive probability used in
the simulation is $0.1\times\text{TP}$, while \id{test\_fp} is read, stored, exposed in the interface and
never applied. The result is that changing the test in the menu changes the false-positive rate only
through its correlation with TP, and always at the fixed 10\% ratio --- so the menu's FP column is
decorative. \emph{(ii)} The aggressiveness cure odds ratio (Section~\ref{sec:agg}) is applied on the
clinical cure path but \emph{not} on the screening cure path, so an aggressive tumour detected by screening
is, in that one respect, treated as if it were not aggressive. Both are small in magnitude and neither
breaks the model, but they mean the model's screening module does not fully honour its own parameter set.}

\subsection{What screen detection does}

\mech{If an eligible, participating agent has preclinical cancer (it has cancer and its onset age has
passed), a positive test marks it screen-detected, records the screening age, and determines the stage at
detection by counting how many crossing ages have already passed at that moment (the same derived-stage
device as at clinical diagnosis, Section~\ref{sec:stage}). A cure is then drawn with the \emph{same} cure
model as clinical diagnosis --- the age-efficiency factor times the stage-specific treatment efficiency ---
but with an independent random draw. If the draw cures, the patient is flagged screening-cured and receives
a screening death age no earlier than five to ten years after detection; if it does not cure, the
clinical-cure flag is \emph{cleared} and the untreated terminal age stands.}

\why{The single most important design decision here is that screen detection re-uses the clinical cure
model with a fresh draw, and that the cure probability depends on the stage \emph{at which the cancer was
found}. That is precisely where screening benefit comes from in this model: earlier detection means an
earlier stage, and an earlier stage means a higher cure probability ($0.7$ at stage I versus $0.2$ at
stage IV). No separate ``screening is better'' effect is assumed --- the benefit is entirely a stage-shift
consequence, which is the honest way to model it and the way that is falsifiable against trial results.}

\why{\emph{The override, and why it must be stated.} Screen detection \emph{replaces} the clinical
trajectory: a failed screening cure clears the clinical-cure flag, so a patient who would have been cured
clinically can end up dying of cancer once detected by screening. The paired lives-saved accounting is
protected against double counting --- an agent is counted as saved only if it is screening-cured \emph{and}
was not clinically cured --- but the override means the model does not treat ``the screen-detected patient
who would have done fine anyway'' as a neutral case. The reported screening effect is therefore the net of
a benefit channel (earlier detection, lower stage, more cures) and a harm channel (overriding a favourable
clinical outcome). In the reference run only 5 patients are screening-cured against 22 screen-detected,
which is consistent with the harm channel being real rather than a corner case.}

\subsection{False positives}

\mech{A participating agent without preclinical cancer tests positive with probability $0.1\times\text{TP}$;
the event is counted per year and contributes to the reported false-positive curve. It has no other
consequence: no follow-up procedure, no work-up cost, no anxiety, no effect on future participation.}

\why{The model counts what it can count. False positives matter to a screening programme mainly through the
burden of unnecessary follow-up, its costs, and the possibility of behavioural change --- and the packaged
inputs contain information about none of these. Counting them keeps the phenomenon visible in the results,
and makes the rate's inertness (previous subsection) obvious to anyone who varies the test menu, without
inventing consequences. Deliberately not modelling a known effect is defensible; silently modelling it with
invented numbers is not.}

\alt{\begin{itemize}[nosep]
  \item \textbf{A continuous-time screening process with exact instants.} Removes the discretisation of
        screening ages and lets detection happen the moment a tumour becomes detectable; requires an
        explicit detectability threshold, which the data do not provide, and it gives up the clean annual
        grid.
  \item \textbf{Correlated (``sticky'') participation, never-attenders, invitation-response
        heterogeneity.} A well-supported and important refinement: screening benefit is concentrated in
        attenders, and a programme with 50\% attendance by a random half behaves differently from one with
        100\% attendance by half the population. Needs individual attendance data; the current design
        cannot express it.
  \item \textbf{Sensitivity and specificity instead of TP/FP.} Mathematically equivalent to what is applied
        (sensitivity $=$ TP, specificity $=1-$FP); the present framing is chosen because the model applies
        the test to an agent known to have preclinical cancer, which is the definition of sensitivity.
  \item \textbf{A size-dependent test sensitivity} (the test improves as the tumour grows). What real tests
        do, and it strengthens the length bias; it requires a link between the Gompertz ``size'' and
        detectability, plus detection-probability-by-size data --- a natural companion to the
        size-threshold alternative of Section~\ref{sec:nonident}.
  \item \textbf{Modelling the downstream consequences of a false positive} (follow-up procedures, their
        costs and complications, anxiety, reduced future attendance). The right long-run direction; it
        needs health-economic and behavioural data that are absent.
  \item \textbf{Route-specific treatment realism.} The present assumption is that a screen-detected cancer
        and a clinically detected cancer of the same stage are treated identically. Real screen-detected
        tumours often have more favourable biology, so this may understate the benefit; the alternatives
        are a route-specific cure parameter or route-specific survival from a trial.
  \item \textbf{An explicit interval-cancer channel} (cancers surfacing between rounds). It arises here
        endogenously, because the diagnosis hazard keeps operating between rounds; making it explicit
        means measuring it, which trial data would allow.
  \item \textbf{Programme scale-up over time} rather than switching on in year 10. A scenario and
        reporting question rather than a structural one.
\end{itemize}}

\section{Death Attribution and the Paired Comparison}
\label{sec:death}

\mech{The model never \emph{draws} a cause of death: it \emph{derives} one, each year, from the ages and
flags it already stores. The categories are: $0$ = no death recorded (alive), $1$ = non-cancer (natural)
death, including a cancer patient whose natural death age arrives first, $2$ = cancer death on the untreated
natural-history path, $3$ = death averted by screening (the patient would have died of cancer but was cured
by the screen), $4$ = death from other causes after a clinical cure.

Two mortality series are produced by the \emph{same} agents in a \emph{single} run: the untreated series,
using each patient's natural-history cancer-death age, and the screening-adjusted series, using the
effective death age (the screening death age for screen-detected, screen-cured patients, and the untreated
age otherwise). Screening benefit is reported as \emph{lives saved} --- patients who are screening-cured,
were not clinically cured, and whose untreated cancer death would have fallen inside the simulated window
(3 people in the reference run) --- and \emph{years of life saved}, the accumulated difference between the
two death ages (14.4 years in the reference run, about 4.8 years per person saved).}

\why{\emph{(i) Deriving the cause makes the categories mutually exclusive and exhaustive by construction.}
Every agent falls into exactly one category in every year, so the counts add up and no separate bookkeeping
can drift out of step.

\emph{(ii) Both arms come from the same agents, so the comparison is paired.} This is the single largest
statistical-efficiency decision in the model. Because the untreated trajectory is already stored for every
patient, the screening effect is a within-individual difference, and the variance of that difference is
orders of magnitude smaller than the variance of a difference between two independently simulated
populations. A screening benefit of a few lives per 20\,000 agents is measurable at all precisely because
of this design; with two independent arms it would be buried in sampling noise.

\emph{(iii) No re-simulation means no seed artefacts in the contrast.} The counterfactual is read from a
stored number rather than obtained by running a control scenario, so the contrast does not inherit the
control run's randomness and cannot be tuned by choosing a seed.

\emph{(iv) The ``not clinically cured'' clause is the correct counterfactual.} An agent counts as saved only
if screening cured it \emph{and} it was not going to be cured anyway. Without that clause, every early
detection would count as a success, including the patients who would have been cured regardless --- which
is exactly the error that makes naive screening evaluations overstate benefit.}

\why{\emph{Three honest qualifications.}
\emph{(a)} Agents leave the population only at their natural death age, so a recorded cancer death does not
remove the agent (\S3.2): the age structure follows the all-cause life table and cause-specific rates use
the general-population denominator. This is the standard cause-specific device under an independence
assumption, and it is a natural target for validation.
\emph{(b)} One branch of the saved accounting tests whether the natural death precedes the untreated
cancer death, i.e.\ whether the person would \emph{not} have died of cancer. It is empty in the reference
run, and it is structurally unable to contribute years (because the screening death age is never earlier
than the untreated one), so it can only inflate the people counter. It is a latent inconsistency rather
than a present error, and it should be tightened.
\emph{(c)} Restricting lives saved to deaths that would have occurred inside the simulated window makes the
reported figure a \emph{within-horizon} number, not a lifetime number: a person saved whose untreated death
would have fallen in year 45 of a 30-year run is not counted. That understates lifetime benefit and makes
the headline figure depend on the horizon --- worth stating explicitly wherever it is quoted.}

\alt{\begin{itemize}[nosep]
  \item \textbf{Two independent populations (a screened and an unscreened cohort).} Simpler to implement and
        easier to explain; the contrast then carries the full Monte Carlo variance, so the same precision
        needs far more agents, and the answer becomes sensitive to the pair of seeds chosen. It also cannot
        express ``this same person's two possible lives''.
  \item \textbf{A same-seed twin run with screening disabled.} Recovers a paired contrast without storing
        the untreated ages, at the cost of doubling the runtime and of relying on the seeds aligning
        exactly. The current design obtains the same information for free.
  \item \textbf{Excess-mortality / relative-survival attribution} (cancer deaths as the excess over a
        background life table, with attributable fractions). Standard in registry analysis and it does not
        need a modelled cure mechanism; it abandons the exact paired accounting, needs cause-deleted
        background mortality, and makes ``saved'' a probability rather than a count.
  \item \textbf{Probability-of-causation attribution} (a death is partly attributable to each cause).
        More faithful to how deaths are actually certified; it destroys the clean integer counts that the
        cost and benefit outputs are built on.
  \item \textbf{Lifetime accounting with survival extrapolation} beyond the horizon. Gives the genuinely
        comparable ``lives saved'' figure used in trial reporting; needs extrapolation beyond the data and
        a discount rate for the years.
  \item \textbf{QALY- or DALY-weighted outcomes} instead of raw lives and life-years. The right currency for
        policy comparison; needs utility weights, a discount rate and a perspective (health system versus
        societal) --- none of which the packaged inputs contain.
\end{itemize}}

\section{Reported Outcomes}

\subsection{Counting, rates and their denominators}

\mech{All outputs are accumulated as counts indexed by \emph{year} and \emph{age}: the age distribution, the
number at risk, incident cases (agents reaching their onset age), diagnoses (agents reaching their diagnosis
age), and cancer deaths in both the untreated and the screening-adjusted series. A rate at age $a$ is then
the total number of events at that age divided by the total person-years at risk at that age, i.e.\ an
average annual rate over the horizon. Two deliberately different incidence-type curves are produced:
\emph{incidence} counts the (unobservable) preclinical onset, and \emph{diagnosis} counts the (observable)
clinical detection.}

\why{Counting on the same grid as the data means the model's output can be compared with the published tables
used to calibrate it without any re-binning, which is what makes calibration and validation audits possible
at all. Splitting the two incidence-type curves is a small design decision with a large diagnostic payoff:
because the onset curve is a modelled quantity and the diagnosis curve is a calibrated one, their divergence
is exactly the assumed lead time made visible. Anyone can therefore see how much of the model's output
depends on the unobservable quantity --- which is the whole point of the lead-time sensitivity analysis.}

\why{\emph{The denominator is the general population, not the cancer-free population.} Because agents are
removed only by natural death (\S3.2), the at-risk denominator for cause-specific rates includes people with
a recorded cancer death. This is the standard cause-specific construction under the independence assumption;
it is also why the reported rates should not be interpreted as ``probability of cancer death among cancer
patients''. Crude rates by age are reported as well, so a user comparing two populations must age-standardise
before drawing conclusions.}

\alt{\begin{itemize}[nosep]
  \item \textbf{Age-standardised rates} (direct standardisation to a reference population) instead of crude
        rates. Necessary whenever two datasets with different age structures are compared --- which is the
        normal case here, since the packaged datasets differ by country. Computing the standardisation in
        the model rather than leaving it to the reader is a small, high-value improvement.
  \item \textbf{Cumulative risk / lifetime risk} (``the probability of a diagnosis before age 75''). The
        headline metric in most cancer statistics, and more interpretable to a non-specialist than an
        annual rate; it is a re-expression of what the model already computes.
  \item \textbf{Incidence-based versus mortality-based outcome framing.} The epidemiological convention is
        to frame results either in terms of cases avoided or deaths averted; the model reports both, which
        is the superset, but the choice matters when a downstream decision model expects one specific
        currency.
  \item \textbf{Interval-based rather than single-year counts} (five-year age bands, calendar periods).
        Aligns with how statistics are published and reduces the sparsity of small counts; loses resolution
        exactly where the model is most informative (the screening window).
  \item \textbf{Rate smoothing} (as is applied to the sensitivity curves). Makes charts readable but changes
        the numbers; it should stay a presentation-layer choice, as it is at present.
\end{itemize}}

\subsection{Survival, stratified by stage and aggressiveness}
\label{sec:survival}

\mech{Two overall cause-specific survival curves are produced --- one on the untreated series and one on the
screening-adjusted series --- both measured from diagnosis. In addition, and more recently, the model
reports cause-specific survival from diagnosis \emph{stratified by stage and by aggressiveness} (each of
stages I--IV plus the pooled curve, in both the aggressive and the non-aggressive subgroup), together with
the \emph{cure fraction} of every one of those strata. The event is defined precisely: a patient dies of
cancer if it is not cured (clinically or by screening) and its untreated cancer-death age is reached before
its natural death age. Cured patients, and uncured patients who die of something else first, are censored.
The curve is accumulated as a discrete hazard, $H_t = H_{t-1} + d_t/n_t$ for event counts $d_t$ and numbers
at risk $n_t$, with $S_t = \exp(-H_t)$ over a 30-year horizon. In the reference run the pooled cure fraction
is $13.1\%$ for aggressive and $27.2\%$ for non-aggressive patients.}

\why{\emph{(i) Cause-specific rather than all-cause} because the quantity of interest is the cancer's effect,
and because all-cause survival mixes in the larger non-cancer mortality --- which differs between the
packaged datasets, so an all-cause curve would make cross-dataset comparison meaningless.

\emph{(ii) Measured from diagnosis} because that is the clinically meaningful time origin and the one used
in registry survival reporting, so the output is directly comparable with published figures.

\emph{(iii) The ``cured $\Rightarrow$ censored'' convention is the correct cause-specific device.} Once a
patient is cured the cancer can no longer kill, so for a cancer-specific curve that patient is no longer at
risk of the event. Counting a cured patient as an event would be plainly wrong; treating it as censored is
the standard construction, and it is the same rule that the death-attribution layer uses (\S10), so the two
cannot disagree.

\emph{(iv) Stratification is the Simpson's-paradox guard.} Because aggressiveness rises steeply with stage
($0.096$ at stage I to $0.639$ at stage IV, \S7), the pooled aggressive-versus-non-aggressive curves are
mostly measuring stage composition. The per-stage comparison is the only one that isolates the odds-ratio
assumption, so reporting pooled curves alone would invite exactly the wrong conclusion.

\emph{(v) Reporting cure fractions alongside the curves makes the central mechanism auditable.} The model's
benefit channel is ``earlier stage $\rightarrow$ higher cure probability'', so the reader should be able to
see how many patients the cure abstraction actually rescues, by stage and by aggressiveness --- and to check
that the reported fractions move in the direction the parameter implies.}

\why{\emph{The caveat that follows from \S8.} These curves are conditional on the assumed cure
probabilities, which are hand-set and confounded with the fitted terminal hazard. A survival curve can
therefore be ``right'' for partly compensating reasons. The honest reading is: the curves are a diagnostic
of the model's internal mechanism and a tool for scenario comparison, not an independent estimate of cancer
survival. Comparing them against externally published stage-specific survival is the test that would settle
it.}

\alt{\begin{itemize}[nosep]
  \item \textbf{All-cause survival.} Unambiguous --- it needs no judgement about the cause of death --- but
        confounded by background mortality, so it cannot be compared across populations with different life
        tables.
  \item \textbf{Relative / net survival} (observed survival divided by expected survival from a background
        life table). The registry standard and the only construction that makes the model's survival
        comparable with published figures; it needs background mortality by age, sex and period.
  \item \textbf{Excess-hazard modelling of survival} (the cancer-attributable hazard as a function of stage,
        age and time since diagnosis). Closer to how the data behave and it removes the need for a binary
        cure; it requires individual follow-up to estimate.
  \item \textbf{A fitted parametric survival model} (Weibull, Gompertz, log-normal) instead of a
        non-parametric curve. Compact summary statistics, extrapolation, and easy comparison of models; it
        adds a shape assumption that the data cannot always support.
  \item \textbf{Conditional survival} (survival given that five years have already been survived). Clinically
        the most interpretable framing of the ``cure plateau'', and it makes the cure abstraction visible
        without adding a parameter.
  \item \textbf{Hazard ratios between the arms} instead of full curves. A compact effect measure that is easy
        to tabulate; it discards the time pattern, which is the aspect screening most affects.
  \item \textbf{Stratification by additional axes} (age band, route of detection, dataset). Cheap to
        compute, and it would expose interactions the pooled reporting hides; the risk is small strata and
        noisy curves, which is why the horizontal resolution should stay at stage and aggressiveness.
\end{itemize}}

\subsection{Stage shift: why the two stage distributions are not comparable}
\label{sec:stageshift}

\mech{The model reports the stage distribution of clinically diagnosed cancers and, separately, of
screen-detected cancers. In the reference run the clinical mix is $0.225/0.250/0.284/0.241$ (mean stage
$2.55$), while the screen-detected mix ($n = 22$) is $0.000/0.318/0.500/0.182$ (mean stage $2.86$). For
those same 22 agents, however, the screening stage \emph{minus} the clinical stage averages $-0.68$: 13 of
22 were detected at an earlier stage (three by two stages, ten by one), eight at the same stage and one
later, with a mean of $4.8$ years between detection and the clinical diagnosis that would otherwise have
occurred.}

\why{These numbers must be read together, because the two stage distributions are \emph{not like-for-like}
and reporting only the marginals would invite the wrong conclusion in either direction. Screening can only
catch a person who happens to be in the preclinical phase at a round, and a tumour with a long sojourn has
more opportunities to be caught. That is length bias: the screen-detected group is enriched with
slowly growing tumours, which have had \emph{more} time to progress. It is therefore entirely possible ---
and it happens in the reference run --- for the screen-detected mix to look worse than the clinical mix in
cross-section, while every within-person comparison shows that the screened tumour was found earlier than
it otherwise would have been.

The model's benefit mechanism is precisely the within-person shift: a lower stage at detection implies a
higher cure probability (Section~\ref{sec:agg}, stage I $0.7$ versus stage IV $0.2$). Measuring the
mechanism therefore requires comparing each detected agent with its own counterfactual, which the model can
do because it stores the clinical stage for every agent. The mean detection-to-clinical-diagnosis interval
--- about $4.8$ years here --- is the quantity that caps the achievable benefit: it is how much earlier the
screening route acts.}

\why{\emph{Sample-size warning.} Twenty-two screen-detected cancers in 20\,000 agents is a small count, so
the screen-detected stage mix carries wide uncertainty and should not be quoted as a point estimate without
an interval or an ensemble over seeds. The within-person mean shift is more stable because it is a paired
quantity, which is the same reason the paired design was chosen for the benefit accounting (\S10). This is
also a practical argument for the standardised, automated sensitivity sweep described in
Section~\ref{sec:sens}: single-run screening statistics are too noisy to inspect one at a time.}

\alt{\begin{itemize}[nosep]
  \item \textbf{Aggregate stage shift over all cancers} (screened and unscreened pooled). Simplest to
        present and the least informative, because it is exactly the comparison that length bias distorts.
  \item \textbf{The paired within-person stage shift as a first-class output} (its mean, median and full
        distribution, plus the detection-to-diagnosis interval). This is the recommendation: it is the
        quantity that actually drives the benefit, it is stable because it is paired, and it cannot be
        misread as a population-level shift.
  \item \textbf{The distribution of the screen lead time} rather than its mean. Shows the tail of very long
        sojourns that dominates the aggregate stage mix.
  \item \textbf{TNM-based shift} (tumour size, nodal involvement, distant metastasis) rather than the
        four-stage summary. The clinically meaningful description, and it would let the model's stage shift
        be compared with trial results; it needs size and nodal data.
  \item \textbf{Age- and stage-standardised incidence and mortality by detection route} (the classical
        stage-shift analysis used in the screening literature). The standard that would make the model's
        output comparable with published screening evaluations; it requires standardisation and a
        like-for-like denominator, both of which are straightforward additions.
  \item \textbf{``Proportion detected at stage I--II'' as the headline screening metric.} Intuitive, easy to
        communicate, and it exposes the length-bias caveat immediately --- which is either a feature or a
        trap depending on how it is worded.
\end{itemize}}

\section{Costs}

\mech{The cost model is three lines of arithmetic. The screening cost is the number of tests performed
times a unit test price. The treatment cost is the number of \emph{cured} patients at each stage times that
stage's unit treatment price, summed over stages ($1,2,3,4$ by stage in the packaged configuration). The
total is their sum. In the reference run: $13\,111$ tests, screening cost $13\,111$, treatment cost $398$,
total $13\,509$.}

\why{The prices are dimensionless because no trustworthy cost data are shipped, and a screening model whose
headline message is a strategy \emph{ranking} needs the relative prices to be defensible rather than the
absolute ones. Encoding ``a test costs one unit, treatment costs one to four units depending on stage''
captures only the ordinal, uncontroversial clinical fact --- later-stage treatment is more expensive --- and
leaves the user free to substitute real tariffs. For the same reason nothing in the calibration touches
costs: they cannot silently absorb a misfit elsewhere.

Keeping the module to three lines is itself a design decision. A richer cost model would have to make
assumptions (discount rate, perspective, follow-up protocols, complications) that the packaged data cannot
support, and those assumptions would then drive the cost comparison --- whereas here the arithmetic can be
checked by hand and the sensitivity of the ranking to the prices is transparent.}

\why{\emph{Three limitations that should be read as caveats on any cost figure.}
\emph{(a)} Treatment is charged only for cured patients, so under a payer's perspective the true cost of a
screening programme is understated (patients who are treated and die are not charged) and the
cost-per-life-saved arithmetic is not a cost-effectiveness ratio in the usual sense.
\emph{(b)} Three of the four price parameters are inert: \id{screening\_price} (a per-person programme
cost), \id{complications\_price} and \id{reoccurrence\_cost} are declared, stored, exposed in the interface
and never used by the simulation. Treating a parameter set as configuration while some of it is decorative
is a trap for a user who changes a price and sees nothing happen.
\emph{(c)} There is no discounting, no time preference, no perspective statement and no cost for a false
positive, even though a false positive is the single most common costly event in a real programme.}

\alt{\begin{itemize}[nosep]
  \item \textbf{A full cost-effectiveness analysis} (cost per life-year or per QALY gained, discounting,
        incremental ratios against a comparator, willingness-to-pay thresholds, and probabilistic
        sensitivity analysis over the prices). This is what a health-economic consumer expects, and the
        model already computes the two ingredients (costs and years of life saved). It requires a
        perspective, a discount rate, utility weights and a comparator strategy --- all of which are
        external decisions rather than data.
  \item \textbf{Treatment-pathway costing} (diagnosis work-up, surgery, chemotherapy, radiotherapy,
        palliative and terminal care, each with a state-specific cost). The standard in decision-analytic
        models, and it makes treatment costs respond to the stage at diagnosis in a clinically meaningful
        way; it needs pathway and tariff data.
  \item \textbf{A health-system perspective with real tariffs,} versus \textbf{a societal perspective}
        adding productivity losses, informal care and travel time. The two can differ by a factor, and the
        choice belongs to the user; the model should simply state which one it assumes (at present,
        neither).
  \item \textbf{Micro-costing of the screening episode} (invitation, kit, laboratory handling, and the
        follow-up procedure for a positive test). The dominant real-world cost driver at the margin, and the
        natural place for the currently inert \id{screening\_price} and for the unmodelled false-positive
        consequences.
  \item \textbf{Charging for everyone treated, not only for those cured.} A one-line change that would make
        the cost figure behave as a payer expects, and would make the screening programme look
        correspondingly less attractive.
  \item \textbf{Cost by year since diagnosis} (the first year of treatment costs far more than subsequent
        years). A well-established pattern in the literature, and it would matter for any comparison that
        shifts diagnoses earlier in time.
\end{itemize}}

\section{Elements That Exist But Do Nothing}
\label{sec:inactive}

\mech{The configuration surface and the interfaces expose a number of elements that have no effect on any
result. They are listed here rather than quietly omitted, because a user who changes a parameter and sees
nothing happen has been misled by the model, not by their own mistake. Every entry below was verified by
searching for actual reads in the simulation path.

\begin{center}
\footnotesize\setlength{\tabcolsep}{3pt}
\begin{tabular}{@{}p{4.3cm}p{3.8cm}p{7.2cm}@{}}
\toprule
\textbf{Element} & \textbf{What it was for} & \textbf{Actual status} \\
\midrule
Recurrence probability & recurrence after cure & Draws an exported flag only
(\id{reoccurrence\_probability} $=0.2$); no effect on survival, death timing or cost. \\
Post-treatment mortality & extra mortality after treatment & Declared and never read
(\id{treatment\_mortality\_adjustment} $=1.0$). \\
Treatment complications & complication rate & Declared and never read
(\id{complications} $=[0.1]$). \\
Complication cost & cost per complication & Declared and never read
(\id{complications\_price} $=1.0$). \\
Recurrence cost & cost per recurrence & Declared and never read (\id{reoccurrence\_cost} $=1.0$). \\
Programme cost per person & per-person screening cost & Declared and never read
(\id{screening\_price} $=1.0$); only the test price is charged. \\
False-positive rate & the test's false-positive rate & Read, stored and displayed but never applied
(\id{test\_fp}, \id{test\_fp\_selected}): the simulated rate is $0.1\times$TP. \\
Risk factors & relative risks for exposures & Declared and never read by the simulation
(\id{factors\_rr}, \id{age\_irrelevant\_factors\_rr}, \id{factor\_age\_correction}). \\
Factor accessors & helpers for those tables & Never called (\id{has\_factor}, \id{get\_factor\_rr}). \\
Legacy configuration & the old key--value format & Reachable only by pointing the model at a \id{.txt}
file (\id{from\_legacy\_txt}, \id{\_csharp\_to\_python\_name}); no packaged configuration uses it. \\
Duplicate cure helper & the age-cure step function & Never called (\id{get\_age\_cure\_efficiency}); the
simulation uses its own copy of the same logic. \\
Risk-to-distribution helper & annual risks $\rightarrow$ death-age distribution & Never called
(\id{risks\_to\_distribution}); the model uses observed deaths directly. \\
Curve-evaluation helpers & evaluating the fitted Gompertz curve & Never called (\id{compute\_output},
\id{compute\_output\_train}, \id{compute\_train\_ages}, \id{gompertz\_v}, \id{compute\_stage\_age}); the fit
and the simulation both use the closed-form crossing times. \\
Vectorised hazard draw & a vectorised first-success draw & Never called (\id{PiecewiseHazard.T\_batch}); the
population generator does this inline. \\
Scalar sampling helpers & single-value sampling & Never called (\id{Distribution.generate\_single},
\id{RandomState.next\_double}, \id{next\_doubles}). \\
\bottomrule
\end{tabular}
\end{center}}

\why{Two principles justify listing this rather than deleting it silently.
First, \emph{an unimplemented mechanism should be visible as an unimplemented mechanism}. Several of these
entries (treatment mortality, complications, recurrence) correspond to real clinical effects that the model
does \emph{not} represent; leaving them in the configuration while they do nothing invites a reader to
assume they are active, and the resulting model is more dangerous than a smaller one that admits its scope.
Second, \emph{the size of the inert surface is a measure of the gap between intent and implementation} ---
here it is a fairly concentrated set of treatment-pathway and recurrence features, which is a useful signal
about where the model's ambitions outran its inputs.}

\why{\emph{The weakness this exposes.} Inert parameters are a usability trap. A user who sets a false-positive
rate, a complication rate or a screening price and observes no change has no way to tell a bug from a
non-feature, and the interface's willingness to display the parameter implies it works. The current code
does carry in-line comments labelling several of them as inactive, which is the right instinct, but a
comment is not a contract.}

\alt{\begin{itemize}[nosep]
  \item \textbf{Delete them.} Cleanest: a parameter set that contains only live parameters cannot mislead.
        The cost is losing the documented intent, which is real information for whoever implements the
        treatment pathway later.
  \item \textbf{Keep them, but label them in the interface as inert} (grey out, or annotate ``not used in
        the current model'', with the reason). Preserves the intent and removes the trap --- this is the
        smallest change that fixes the actual problem.
  \item \textbf{Fail loudly when an inert parameter is set to a non-default value} (a validation warning at
        load time). Guarantees that nobody believes a setting they just made is in force.
  \item \textbf{Implement the treatment pathway}, turning recurrence, complications and treatment mortality
        into active mechanisms. The correct long-run answer, and the only one that makes the parameters
        meaningful; it needs treatment and complication data.
  \item \textbf{Move the legacy loader out of the main path} (or delete it) so that a single configuration
        format remains. A second parsing path is a second surface for silent divergence --- the two loaders
        would have to agree on the meaning of every field, and nothing tests that they do.
\end{itemize}}

\section{Calibration}
\label{sec:calib}

\subsection{Three separate likelihoods, and everything else by hand}

\mech{When calibration is requested, the model solves three independent fitting problems.

\begin{enumerate}[nosep]
  \item \textbf{Diagnosis hazard} (6 coefficients): a Poisson likelihood on the observed age-specific case
        counts, $C(a)\sim\mathrm{Poisson}\bigl(h(a)P(a)\bigr)$, where $h(a)$ is the piecewise log-hazard of
        \S4. Minimised with L-BFGS-B under $\pm5$ bounds; it converges in a fraction of a second.
  \item \textbf{Gompertz shape and rate} (3 parameters $K$, $C$, $B_{\text{pop}}$): least squares on the
        three crossing times \emph{derived from the stage mix} (\S6.3), plus penalties that keep the
        carrying capacity above the terminal stage, keep the rate positive, and mildly regularise the
        capacity so the fit cannot run into a degenerate huge-capacity/tiny-rate corner. In the reference
        fit the residual is of order $10^{-3}$ --- the targets are reproduced essentially exactly.
  \item \textbf{Terminal-phase hazard} (1 parameter $\lambda$): a Poisson likelihood on the observed
        age-specific cancer deaths, $D(a)\sim\mathrm{Poisson}\bigl(e^{\lambda}C(a)\bigr)$ with
        $C(a)=I(a)P(a)$. The reference fit gives $\lambda=-1.05$, i.e.\ a rate of $0.35$ per year.
\end{enumerate}

Everything else is a configured value rather than an estimate: the lead-time mean, $B_{\text{std}}$, the
stage treatment efficiencies, the age cure constants, the aggressiveness odds ratio, the screening schedule,
participation, test accuracy and the prices. The stage-mix inversion that produces the crossing times is
arithmetic, not optimisation, and it runs before the Gompertz fit.}

\why{\emph{(i) Splitting the problem is what keeps assumptions assumptions and estimates estimates.} A
single joint fit over all parameters would be non-identifiable along the two flat ridges of \S5.2 and \S8.2:
the optimiser could trade lead time against growth rate, or cure probability against terminal hazard, and
improve the aggregate fit while moving the parameters arbitrarily far. Splitting the fit into three
well-posed pieces means each parameter is pinned by a target that can actually pin it --- and the quantities
that no target can pin are \emph{not} fitted at all: they are configured and then varied in the sensitivity
analysis. That distinction, enforced by the structure of the code rather than by discipline, is the
guardrail that matters most.

\emph{(ii) Each objective is the right likelihood for its data.} Counts of cases and deaths are Poisson
counts, not binomial proportions, and the discrete hazard of \S4 is the natural link. Using the correct
likelihood is not pedantry: the previous formulation used a binomial likelihood on proportions and --- far
worse --- calibrated mortality against the \emph{diagnosis} hazard rather than the cancer-death hazard, which
made the mortality calibration meaningless while still looking plausible.

\emph{(iii) Small and smooth objectives are fast, reproducible and auditable.} Six parameters and a smooth
likelihood converge deterministically from a fixed starting point; one parameter is a line search. Nothing is
fitted by simulation, so calibration takes milliseconds, results are exactly reproducible, and each
likelihood value can be printed and checked.

\emph{(iv) Structural constraints are expressed as smooth penalties.} Every threshold crossing must be
reachable ($\exp K > s$) and the rate must be positive; encoding these as large quadratic penalties keeps the
optimiser out of the degenerate region without a constrained solver or a rejection loop.}

\why{\emph{Three methodological weaknesses to state.}
\emph{(a)} The mortality likelihood is an \emph{approximation}: it treats every case as subject to a constant
death hazard, ignoring that only uncured patients die and that death follows diagnosis with a delay. The
fitted $\lambda$ therefore absorbs whatever the assumed cure probabilities get wrong --- which is exactly the
confound of \S8.2, seen from the calibration side.
\emph{(b)} There is no train/validation split: the model is calibrated and judged on the same data. This is
the weakest methodological point in the whole design, and the remedy is procedural rather than
computational.
\emph{(c)} Not everything that \emph{could} be fitted is fitted. The stage-mix calibration is
dataset-specific by construction (the crossing times are derived from that dataset's stage mix), so the only
components transferable between datasets are the diagnosis hazard and the mortality hazard. A genuine
external validation would therefore have to hold out one of those two, or hold out another population and
accept that the staging calibration is re-derived --- which limits what ``validation'' can mean here, and
should be said wherever a validation claim is made.}

\alt{\begin{itemize}[nosep]
  \item \textbf{Joint Bayesian calibration with priors and posterior predictive checks} (MCMC or ABC). The
        statistically correct way to handle the flat ridges: the assumptions enter as priors and their
        consequences arrive as credible intervals instead of as a set of re-runs. Costs hours instead of
        milliseconds, and the headline numbers then depend on priors that are themselves not derived from
        data.
  \item \textbf{Simulated maximum likelihood} (computing the likelihood from simulation rather than
        analytically). Removes the mortality approximation and would make the cure-versus-terminal-hazard
        ridge explicit rather than hidden; the objective becomes noisy and far more expensive, which in turn
        requires many replications or a stochastic optimiser.
  \item \textbf{Approximate Bayesian computation or history matching.} Likelihood-free and well suited to a
        simulator whose analytic likelihood is only approximate; needs thousands of runs and an explicit
        tolerance, so it is a research workflow rather than an interactive one.
  \item \textbf{A surrogate model plus global optimisation} (the differential-evolution and Optuna paths
        already present in the code). Useful when the objective is multi-modal or noisy; it adds a
        hyper-parameter --- the surrogate's own fit --- to the problem.
  \item \textbf{Manual / expert (Delphi) calibration}, as practised for several large cancer models.
        Transparent, and it can incorporate judgement the data do not carry; it is not reproducible and it is
        hard to audit.
  \item \textbf{Iterated sequential calibration to convergence} (fit the hazards, hold them, fit the staging,
        then revisit the hazards). Makes the coupling between components explicit; it needs a convergence
        criterion, and it can drift if the components are only weakly coupled.
  \item \textbf{A genuine hold-out protocol:} calibrate on one population or period and validate on another
        --- the packaged USA and Germany datasets make this a one-line experiment for the transferable
        components --- split temporally where multiple periods exist, and report the validation error as the
        headline quality metric. This is the highest-value methodological addition available without new
        data.
\end{itemize}}

\section{Randomness and Reproducibility}

\subsection{Two layers of randomness}

\mech{The model uses two layers. \emph{Population-level} draws (initial ages, natural death ages, the
per-agent stream seeds) come from a single PCG64 generator seeded once per run and stored with the results.
\emph{Agent-level} draws inside the cancer-history kernel do not use that generator at all: each agent has a
64-bit stream identifier, and each quantity is drawn as $\id{uniform}(\text{seed}_i + k)$ for successive
small offsets $k$, using a SplitMix64-style avalanche hash. Gaussian draws are Box--Muller on two uniforms
from the same stream; Bernoulli draws are uniform comparisons. The generator state is \emph{thread-local},
so concurrent runs in one process cannot re-seed each other, and the seed is logged with every run.}

\why{\emph{(i) Per-agent streams decouple agents from each other and from the order of work.} An agent's
draws depend only on its own identifier, so the results are invariant to how the loop is laid out, to
vectorisation, and to any future parallelism. That is what allows the kernel to iterate agents in whatever
order is convenient and to be compiled as a tight loop.

\emph{(ii) Neighbouring indices must decorrelate, because the kernel reads them consecutively.} Each agent
draws many quantities from nearby offsets, so the map from index to value has to look like independent noise
for offsets of 1, 2, 3\ldots. It did not. The previous linear congruential generator followed by a single
xorshift step gave $\mathrm{corr}\bigl(u(s),u(s{+}1)\bigr)\approx+0.36$ and
$P\bigl(u(s{+}1)<u(s)\bigr)\approx0.12$ --- instead of $0$ and $0.5$ --- which silently coupled the stage
draw to the lead time, the cure draw, the cancer-death age and the aggressiveness label. Aggressive tumours,
for example, were drawn almost only at stage IV. This was a real modelling error, not an aesthetic one: it
changed the joint distribution of the model's own variables. The avalanche hash fixes it, and a guard test
now pins the invariant ($|\mathrm{corr}|<0.015$ and $|P-\tfrac12|<0.015$ at ten lags).

\emph{(iii) Thread-local state is a correctness requirement for a web service.} The application serves
concurrent runs from separate threads; a process-wide singleton let one run re-seed the stream another run
was drawing from, which broke reproducibility only intermittently and therefore looked like flakiness rather
than a bug.

\emph{(iv) PCG64 for the population-level draws} is a modern, small-state, well-tested generator and it is
the current NumPy default --- the safe choice for quantities that are drawn once per agent in bulk.

\emph{(v) Logging the seed makes every result reproducible in one command}, which is a prerequisite for the
calibration, debugging and validation workflows, and for reproducing a figure quoted in a report.}

\why{\emph{Two honest qualifications.} First, the agent-level streams are \emph{pseudo}-independent by
construction: they are not formally independent, and their adequacy rests on the avalanche property being
verified (which the guard test does) rather than on a theorem. Second, exactly reproducing a run requires
both layers to be seeded consistently, so the seed reported with a result must be understood to control the
agent-level seeds too --- it does, because they are themselves drawn from the seeded population-level
generator, but the dependency is worth stating.}

\alt{\begin{itemize}[nosep]
  \item \textbf{A single sequential RNG stream for everything.} Simplest possible design; it makes every
        result depend on the exact order of operations, which forecloses vectorisation and parallelism and
        means any code rearrangement changes all outputs. It is also what makes accidental correlations
        across variables easy to create and hard to notice.
  \item \textbf{A full generator object per agent} (each with its own PCG64 state). Statistically excellent
        and trivially independent; constructing one object per agent costs orders of magnitude more than
        hashing two 64-bit words, and the kernel is executed for every agent in the population.
  \item \textbf{Counter-based generators (Philox, Threefry).} Built precisely for this pattern --- a counter
        plus a key streams out values --- with documented statistical quality and cheap parallel jumps; the
        natural upgrade path, and it would remove the reliance on a hash's avalanche behaviour.
  \item \textbf{Common random numbers across scenarios} (sharing the random shocks between the screened and
        unscreened arms to shrink the variance of the difference still further). The paired design of \S10
        already captures most of this benefit by storing the counterfactual rather than re-simulating it.
  \item \textbf{Replicate runs with reported intervals} rather than a single run (seed averaging). The honest
        way to present a stochastic model's output, especially for tail quantities such as lives saved, where
        a single run's count can move by several units; it multiplies the runtime by the number of
        replicates, which the interactive interface cannot currently afford.
  \item \textbf{Antithetic or quasi-random (Sobol) sampling} for the population-level draws. Reduces
        variance for integration-like quantities; much less useful for rare tail events, which is what the
        interesting outputs are.
  \item \textbf{A recorded draw-by-draw random-number log} for bit-level reproducibility under any future
        code change. Overkill today, but it is the standard a formally certified model is eventually held
        to.
\end{itemize}}

\section{Sensitivity Analysis as the Uncertainty Statement}
\label{sec:sens}

\mech{The sensitivity analysis re-runs the whole model for multiples of the lead-time mean --- $0.5$,
$0.75$, $1.0$, $1.25$ and $1.5$ by default --- and reports the incidence, mortality, screening-mortality,
survival and years-saved curves plus the stage distributions for every factor side by side. For a factor
$f$ the lead-time mean becomes $f\mu$ \emph{and} the stage-crossing times are rescaled by the same factor,
after which the Gompertz is refitted to the rescaled targets, so the model still reproduces the same stage
mix at the new time scale. All other parameters stay at their fitted reference values, and the random seed
is the same for every factor. Display smoothing is applied to the curves, never to the stored numbers.}

\why{\emph{(i) The honest response to an unidentifiable parameter is a range, not a value.} The lead time
cannot be estimated from the available data (\S5.2), so the only defensible thing to do with it is to show
what the answer looks like across its plausible range. The sensitivity analysis \emph{is} the model's
uncertainty statement about its one structural assumption --- not an optional extra.

\emph{(ii) Rescaling the crossing times together with the lead time is what keeps the comparison
meaningful.} Because $t_s=-\mu\ln P(\text{stage}\ge s)$, changing $\mu$ while holding the crossing times
fixed would change the stage mix and thereby confound the time-scale question with a staging change. Moving
both together holds the observable stage mix constant, so the curves differ \emph{because} of the time-scale
assumption and nothing else.

\emph{(iii) A common seed across factors is a variance-reduction device, and here it is essential.} Sharing
the random shocks means the between-factor differences are the model's genuine response to the assumption,
with the Monte-Carlo noise largely cancelled out of the comparison. Without it, and with as few screening
events as the reference run produces, the spread between curves could be dominated by sampling variation
--- a user would be reading noise as sensitivity.

\emph{(iv) Smoothing is presentation-only} so that the numbers reported and exported remain exact, and the
parallel execution path exists so that five full re-runs stay interactive, with an automatic threshold that
avoids paying process start-up for jobs too small to benefit.}

\why{\emph{Two limitations that a reader of these plots should know.} First, \emph{only the lead time
varies}. The other hand-set quantities --- $B_{\text{std}}$, the cure probabilities, the aggressiveness odds
ratio, participation, test accuracy --- are equally unverified, and a set of curves that varies only one axis
can easily be read as ``this is the only soft assumption''. Second, the analysis reports \emph{curves} rather
than a decision metric, so the quantity a decision-maker would actually want ranged (lives saved, or cost
per life saved, under each factor) is not tabulated and has to be assembled from the metrics block by hand.}

\alt{\begin{itemize}[nosep]
  \item \textbf{A standardised sweep over every soft parameter}, with the results tabulated as decision
        metrics rather than as families of curves. One-at-a-time analysis is a lower bound on uncertainty
        (it cannot show interactions); it is easy to run because the model is fast, and it is the most
        valuable single addition available.
  \item \textbf{Global sensitivity analysis} (variance decomposition, Sobol indices, elementary-effects
        screening). Apportions the output variance across the assumptions and exposes interactions; needs a
        designed experiment over hundreds or thousands of runs, which is within reach of the current model
        but not of the current interface.
  \item \textbf{Probabilistic sensitivity analysis} (distributions on the soft parameters, Monte Carlo over
        draws, and intervals on the outcomes). The health-economic standard, and it converts
        ``here are five curves'' into ``here is the distribution of the answer''; it requires defensible
        distributions, which for the hand-set parameters means priors.
  \item \textbf{Scenario analysis} (``optimistic test'', ``pessimistic participation'', ``slow tumours'')
        instead of a numeric sweep. More communicable to non-modellers and it forces the assumptions to be
        named; it is a presentation of the same information and loses the systematic coverage.
  \item \textbf{Value-of-information analysis} (which assumption, if researched, would most reduce decision
        uncertainty). The natural next question after a sensitivity analysis, and the one that turns the
        exercise into a research agenda; it needs a decision rule and a loss function.
  \item \textbf{Structural sensitivity:} swapping the mechanism itself (a size-threshold detection model
        instead of a lead-time model, per-stage dwell times instead of crossing times, an onset-driven
        natural history instead of a diagnosis-driven one) and comparing the outputs. No parameter sweep can
        reach this class of uncertainty, and it is where the largest unquantified risk lies.
\end{itemize}}

\section{Assumptions, Direction of Bias, and Falsifiability}
\label{sec:assumptions}

\subsection{The assumption register}

\mech{Every simplification in the model is listed below with the direction in which it pushes the results.
Where the direction is stated, it is the direction relative to a ``true'' model of the same question; where
it is not determinable from the design alone, that is said instead of guessed.

\begin{center}
\footnotesize
\begin{tabular}{@{}p{8.0cm}p{7.6cm}@{}}
\toprule
\textbf{Assumption} & \textbf{Direction of the bias if it is wrong} \\
\midrule
No pre-cancer stage (no adenoma--carcinoma sequence; every cancer arises invasive). & Overstates the
screening benefit: removing a pre-cancerous lesion is an additional protective mechanism the model lacks. \\
No non-progressive lesions. & Overstates the benefit: every screen-detected cancer is assumed to be one
that would have surfaced and could have killed, so overdiagnosis cannot be represented. \\
Diagnosis-driven natural history (no undiagnosed cancer deaths). & Understates the population's cancer
burden; harmless for the screening contrast, misleading for a burden estimate. \\
The absolute preclinical time scale (lead-time mean) is assumed, not estimated. & Not determinable --- this
is precisely why it is swept (\S16). It is the model's largest single assumption. \\
Growth heterogeneity ($B_{\text{std}}$) is set by hand. & Not determinable; too little makes the stage shift
too clean and too favourable, too much does the reverse. Currently not swept. \\
Cure probabilities are set by hand and are confounded with the fitted terminal hazard. & The aggregate
cancer mortality is matched either way; what is not identified is the composition --- how many patients the
``cure'' actually rescues. \\
Aggressiveness does not affect growth. & Understates the biology; the observed stage gradient in survival
is then partly a model construction rather than a mechanism. \\
One primary cancer per person; no second primaries or metachronous tumours. & Slightly understates incidence
and mortality. \\
No treatment-related mortality, no complications. & Understates the harm side of both treatment and
screening (the follow-up after a positive test). \\
No behavioural effects (a false positive does not change attendance; screening does not change lifestyle).
& Not determinable; real programmes show both effects, in opposite directions. \\
Uniform treatment at a given stage, regardless of comorbidity, access or era. & Understates variation and
inequality; the average may be biased in either direction depending on the population. \\
Closed population: no births, no migration. & Only matters over horizons long enough for the age structure
to drift, which 30 years largely is not. \\
Annual discrete time; within-year ordering ignored. & Small, but systematic for the fastest terminal
phases (mean terminal survival is only about three years). \\
Cancer death does not remove the agent from the simulated population; cause-specific rates use the
general-population denominator. & Treats the two causes as independent; small at these rates, but it
understates the correction that a competing-risk denominator would apply. \\
Age is the only exposure dimension: no sex, no region, no time trend, no risk factors. & The model
represents one population in one period; applying it to trends, to another sex, or to a risk-stratified
programme would be invalid rather than merely imprecise. \\
Benefit is measured within the simulation horizon. & Understates lifetime benefit, and makes the headline
figure depend on the horizon. \\
Test performance is a constant (true-, false-positive) pair. & Omits size dependence and all false-positive
harm, so the net benefit is slightly overstated. \\
Screen-detected and clinically detected cancers of the same stage are treated identically. & Ambiguous:
screen-detected tumours often have better biology (understating the benefit), while the override of the
clinical cure path can work the other way (\S9.3). \\
No discounting and no time preference. & Overstates long-run benefits relative to costs when compared with
a discounted analysis. \\
Risk factors, complications, recurrence and treatment mortality exist as parameters but do nothing
(\S13). & No effect; the risk is purely that a reader assumes they are active. \\
\bottomrule
\end{tabular}
\end{center}}

\subsection{What would show the model to be wrong}

\mech{Three classes of check are available, and they differ enormously in strength.

\emph{(i) Internal consistency checks} --- quantities the model is supposed to reproduce from data it was
given: the age-specific incidence curve (fitted), the age-specific cancer mortality curve (fitted, but with
the cure-versus-terminal-hazard confound of \S8.2), and the stage mix at diagnosis. The third is the
informative one, because only its three crossing-time \emph{targets} are taken from the data: the stage mix
itself is then produced by the interaction of the growth model and the lead time. In the reference run it is
reproduced to within a few thousandths ($0.225/0.250/0.284/0.241$ against $0.220/0.256/0.308/0.218$). A
failure here would point squarely at the growth model, the crossing-time inversion, or the per-agent draws.

\emph{(ii) External checks} --- the ones that can actually falsify the model: a screening trial of the same
test in the same age range, where the model's benefit must agree in both sign and magnitude; the observed
mix and rate of interval cancers in a real programme; the within-person stage shift measured in a screened
cohort (\S11.3); and the sojourn time inferred from screen-detected versus interval cancers, which would
identify the one quantity the model has to assume. None of these are possible with the packaged aggregate
data, and all of them are possible with trial or registry screening data.

\emph{(iii) Structural checks} --- consistency of the mechanism with what is known about the disease. The
model has no pre-cancer stage and no non-progressive lesions, so if it predicts a benefit where a trial
found none, those two omissions are the first suspects, in that order. A structural check can also be run as
an experiment rather than as an argument: implement the omitted mechanism and see whether the conclusion
survives.}

\why{Falsifiability cannot be bolted on at the end: it has to be specified before the numbers are seen, or
it degenerates into post-hoc rationalisation. The model is unusually well placed for this because it is fast
and deterministic given a seed --- a criterion can be checked in seconds --- yet it currently has no
pre-registered success criteria and no formal validation against anything outside its own calibration data.
That is the gap between a working model and a defensible one.}

\alt{\begin{itemize}[nosep]
  \item \textbf{Pre-registered success criteria and analysis plan} (which outputs, which tolerance, which
        dataset, decided in advance). Cheap, purely procedural, and it converts validation from a narrative
        into a test.
  \item \textbf{A permanent benchmark suite} (recorded reference outputs under fixed seeds, checked on every
        change). The model already does this for the random generator; extending the same idea to the whole
        simulation would catch unintended changes --- the ones that produce plausible-looking but different
        numbers, which are the dangerous ones.
  \item \textbf{Periodic recalibration against newly published data}, with the previous version's
        predictions compared against the new observations and the discrepancy reported. Turns an annual data
        release into a published back-test.
  \item \textbf{An independent implementation of one component} (for example an independent recomputation of
        the crossing-time inversion and the Gompertz fit) to catch a shared misconception between the
        specification and the code.
  \item \textbf{Decision-analytic validation} (does the model change a decision that would otherwise have
        been made?). The most relevant test for a screening model, and the only one that directly addresses
        whether the modelling error matters.
\end{itemize}}

\section{Summary: Every Element, Its Choice, Its Alternatives}

The table below is the whole document in one place. For each modelling element it states the design that was
chosen, the alternatives that were considered, and what would change upstream if the design were switched.

{\footnotesize\setlength{\tabcolsep}{4pt}
\begin{longtable}{@{}p{2.3cm}p{4.5cm}p{4.0cm}p{4.0cm}@{}}
\toprule
\textbf{Element} & \textbf{Chosen design} & \textbf{Main alternatives} & \textbf{What changes upstream if switched} \\
\midrule
\endfirsthead
\toprule
\textbf{Element} & \textbf{Chosen design} & \textbf{Main alternatives} & \textbf{What changes upstream if switched} \\
\midrule
\endhead
\bottomrule
\endlastfoot
Model class & Agent-based microsimulation; both arms from the same agents & Cohort Markov, age-structured PDE,
renewal equation, semi-Markov, compartmental ODE & The paired lives-saved accounting (the model's statistical
efficiency) is lost; heterogeneity must be discretised into states. \\
Time & Annual discrete steps on an integer age grid & Continuous time with exponential clocks, monthly
steps, event-driven simulation & Screening instants and terminal-phase timing become exact; the calibration
targets stay annual, so the extra resolution is not controlled by data. \\
Population frame & Closed cross-sectional cohort, 30-year horizon & Birth cohort from age 0, open population
with demographic projection, steady state & A birth cohort needs a century-long burn-in; an open population
contaminates the screening contrast with demographic drift. \\
Initial ages & Observed age structure used as a probability mass & Uniform or parametric age profile, age
bands, census microdata & Parametric profiles impose a shape; a uniform start destroys the dominant driver of
incidence. \\
Non-cancer mortality & One sampled natural death age per agent, from observed deaths & Annual Bernoulli
survival draws, parametric mortality law, cause-deleted life tables, frailty & Cause-deleted tables would
separate the two causes properly (needs a cause-specific death column); frailty needs an unidentifiable
distribution. \\
Onset hazard & Piecewise-constant log-hazard on age; the first Bernoulli success defines the diagnosis age &
Smooth parametric hazard, splines, observed rates used directly, mechanistic carcinogenesis model & Splines
are the natural refinement; a mechanistic model would make risk factors modelable but is not identifiable
from incidence alone. \\
What counts as cancer & Diagnosed before natural death; agents past 98 excluded & Onset-driven natural
history, a non-progressive lesion fraction, no exclusion at all & Onset-driven natural history would make
overdiagnosis representable, and is the largest structural upgrade available. \\
Lead time & Exponential with an assumed mean (3.5 years) & Gamma/Weibull sojourn, deterministic lead time,
per-stage dwell times, size-threshold detectability, estimation from trial data & Only trial or registry
screening data identifies the scale; size-threshold detectability would unify detection with the growth
curve. \\
Time-scale identification & Fix the scale; derive crossing times from the stage mix; sweep the scale &
Bayesian prior and posterior, profile likelihood, fix the growth rate instead, two-parameter sojourn & A
posterior converts one number plus a caveat into intervals on every output, at the cost of hours and of a
prior. \\
Growth curve & Reduced Gompertz, $V(t)=\exp(K-\exp(C-Bt))$ & Exponential, logistic, von Bertalanffy,
Weibull/power law, mechanistic ODE, branching process, no growth model & A different curve changes the
closed-form crossing times, and with them the entire staging calibration and the identifiability argument. \\
Growth heterogeneity & Log-normal $B_{\text{ind}}$ with hand-set $B_{\text{std}}=0.1$ & No heterogeneity,
Gamma growth rate, discrete growth classes, random effects on covariates & $B_{\text{std}}=0$ removes length
bias and makes the stage a deterministic function of the lead time. \\
Staging & Derived: the number of size thresholds crossed before diagnosis & Draw from the observed mix,
ordinal regression on covariates, latent multi-state progression, TNM & Drawing the mix makes calibration a
tautology; TNM would need size and nodal data. \\
Aggressiveness & Per-stage Bernoulli label modifying cure odds (OR 0.667) & Growth modifier, continuous
latent score, mixture cure model, subtype-specific survival, no label & Making it a growth modifier
double-counts against $B_{\text{std}}$ and breaks the exact stage-mix fit. \\
Cure & Stage efficiency $\times$ age step function, capped, Bernoulli; cure means the cancer no longer
decides death & Relative-survival tables used directly, mixture cure model estimated from follow-up,
excess-hazard modelling, treatment pathways & Registry survival tables would remove the mechanism and the
ability to ask why a patient survives. \\
Screening schedule & Age window plus fixed period; independent per-round participation & Continuous-time
screening, correlated/sticky participation, never-attenders, programme scale-up & Sticky participation would
make the benefit depend on reaching under-screened groups; the present design cannot express it. \\
Test performance & A (TP, FP) menu of which only TP is applied (FP $=0.1\times$TP) & Sensitivity/specificity
framing, size-dependent sensitivity, explicit follow-up costs & Making FP live is a bug fix;
size-dependent sensitivity is what makes the length bias quantitative. \\
Screen detection & Re-uses the clinical cure model with the stage at detection; overrides the clinical path &
Additive effect, route-specific cure parameter, route-specific survival from a trial & The override creates a
harm channel; an additive design would remove it but understate the benefit. \\
False positives & Counted only, with no consequences & Follow-up procedures, costs, anxiety, effects on
attendance & Modelling them is the correct direction, and it needs health-economic and behavioural data. \\
Death attribution & Derived from the stored ages; five exhaustive categories & Two-population comparison,
same-seed twin run, relative survival, probability of causation & A two-population comparison loses the
pairing and needs far more agents for the same precision. \\
Rates and outcomes & Counts by year and age; crude rates over person-years & Age standardisation, cumulative
risk, interval bands, smoothing & Standardisation is necessary for cross-dataset comparison, and it is
currently left to the reader. \\
Survival & Cause-specific from diagnosis, stratified by stage and aggressiveness, cured patients censored &
All-cause survival, relative/net survival, excess hazard, parametric models, conditional survival & Relative
survival is the registry standard and the only like-for-like comparison with published figures. \\
Stage shift & Both distributions reported; the within-person shift is the mechanism & Aggregate shift only, a
paired shift as a first-class output, TNM shift, standardised stage-shift analysis & Reporting the paired
shift would prevent the length-bias misreading that the marginals invite. \\
Costs & Tests $\times$ unit price plus cured patients $\times$ stage price; dimensionless & Full
cost-effectiveness analysis, treatment-pathway costing, real tariffs, societal perspective, micro-costing &
Charging everyone treated and adding follow-up costs would change the economics qualitatively, not merely in
level. \\
Inactive elements & Documented rather than removed (\S13) & Delete them, label them as inert in the interface,
warn on non-default values, implement the treatment pathway & Implementing recurrence, complications and
treatment mortality makes those parameters meaningful, and adds a real harm channel. \\
Randomness & PCG64 at population level, SplitMix64 hash per agent, thread-local state & A single sequential
stream, a generator per agent, counter-based generators, common random numbers, replicates with intervals &
Replicates with reported intervals are the honest presentation for tail quantities such as lives saved. \\
Calibration & Three separate likelihoods; everything unidentifiable is configured by hand & Joint Bayesian
calibration, simulated likelihood, ABC/history matching, surrogate plus global optimisation, Delphi &
A joint Bayesian fit would quantify the ridges, but would make the headline numbers prior-dependent. \\
Validation & None formal; the stage mix is the only informative internal check & Pre-registered criteria,
hold-out population, benchmark suite, periodic back-testing, decision-analytic validation & Pre-registration
is procedural and free, and it is the largest gap between this model and a defensible one. \\
Sensitivity & One axis (the lead-time scale), common seed, curves reported & A sweep over all soft
parameters, global variance decomposition, probabilistic analysis, scenarios, value of information,
structural swaps & Sweeping all soft parameters would show that the lead time is not the only unverified
assumption. \\
\end{longtable}
}

\section{Where to Read More}

This document is the ``why'' of the model. Its two siblings cover the other angles:

\begin{itemize}[nosep]
  \item \id{model\_for\_dummies.pdf} --- the same model in plain language, for a reader who wants to know
        what the software does rather than why it was built this way.
  \item \id{python\_port\_analysis.pdf} --- the technical and statistical review: state variables, the
        statistical corrections applied during the port, performance, and the methodological improvements
        that would be needed for a formal policy-grade model.
\end{itemize}

The parameter file (\id{config/parameters.toml}) carries a comment against every hand-set parameter saying
what it is and whether it is estimated; the sections above are the long form of those comments. The two
structural assumptions to keep in view while reading any result are the assumed preclinical time scale and
the assumed cure probabilities, because those are the two quantities that the data cannot pin down --- and
they are also the two that the sensitivity analysis, in its present one-axis form, does not fully cover.

\end{document}
